\documentclass[10pt,a4paper]{article}

% ----------------- Paquetes -------------------%
\usepackage[T1]{fontenc}
\usepackage[utf8]{inputenc}
\usepackage[a4paper,margin=2.5cm]{geometry}
\usepackage{mathptmx}             
\usepackage{changepage} 
\usepackage{setspace}
\usepackage[spanish]{babel}
\usepackage[labelfont=bf,labelsep=space]{caption}
\usepackage{amssymb}
\usepackage{amsmath}
\usepackage{ebgaramond}
\usepackage{placeins}
\usepackage{etoolbox}
\usepackage{setspace}
\usepackage{parskip} 
\usepackage{tcolorbox}
\usepackage{needspace}
\usepackage{xparse}
\usepackage{ocgx2}
\usepackage{hyperref}
\usepackage{hypcap}


%%%%%%%%%%%%%%%%%%%%%%%%%%%%%%%%%%%%%%%%%%%%%%%%%%%%%%%%%%%%%%%%
% --- Fija primero tus valores globales "buenos" ---
\setstretch{1.2}                 % Interlineado global
\setlength{\parskip}{0.7\baselineskip}
\setlength{\parindent}{0pt}
\newlength{\GlobalParskip}
\setlength{\GlobalParskip}{\parskip}

\newcounter{eqnum}[subsection]
\renewcommand{\theeqnum}{\arabic{eqnum}}
\newcounter{alphaitem}
\newcounter{numitem}

%% ------------------- Recuadros----------------------%%
\newcommand{\puntosclave}[1]{%
  \begin{tcolorbox}[
    colframe=black,
    title={Puntos clave},
    before upper={\setlength{\parindent}{0.5cm}\setlength{\parskip}{0.5\baselineskip}}
  ]
  #1
  \end{tcolorbox}
}

\newcommand{\modificaciones}[1]{
  \begin{tcolorbox}[
    colback=white,
    colframe=red,
    title={Modificaciones a realizar},
    before upper={\setlength{\parindent}{0.5cm}\setlength{\parskip}{0.5\baselineskip}}
  ]
  #1
  \end{tcolorbox}
}

%% ------------------- Preguntas TEST----------------------%%
% Opciones: radio (single-choice)
\NewDocumentCommand{\OptRadio}{m m m}{%
  \noindent\CheckBox[radio,name=#1,value=#2,width=1.2ex,height=1.2ex]{}%
  \;\textbf{#2.}~#3\par
}

% Opciones: checkbox (multi-choice)
\NewDocumentCommand{\OptCheck}{m m}{%
  \noindent\CheckBox[name=#1,width=1.2ex,height=1.2ex,bordercolor={0 0 0}]{}%
  \;#2\par
}

% Pregunta single-choice (radio)
% Args: 1=id, 2=enunciado, 3=opciones, 4=solución+justificación, 5=imagen (o {}), 6=origen
\NewDocumentCommand{\PreguntaTestSingle}{m m m m m m}{%
  \par\medskip
  \noindent\textbf{#1.}~#2\par\smallskip

    % Imagen (si no es {})
    \IfBlankTF{#5}{}{%
      \begin{center}
        \includegraphics[width=0.75\linewidth]{#5}
      \end{center}
    }

    \begin{Form}
      #3
    \end{Form}

    \par\smallskip
    \switchocg{sol-#1}{\fbox{Mostrar / ocultar solución}}%
    \hfill{\footnotesize #6}\par\smallskip

    \begin{ocg}{Solución}{sol-#1}{0}
      \noindent #4
    \end{ocg}
  \par\medskip
}

% Pregunta multi-choice (checkboxes)
\NewDocumentCommand{\PreguntaTestMulti}{m m m m m m}{%
  \par\medskip
  \noindent\textbf{#1.}~#2\par\smallskip

    % Imagen (si no es {})
    \IfBlankTF{#5}{}{%
      \begin{center}
        \includegraphics[width=0.75\linewidth]{#5}
      \end{center}
    }
    \begin{Form}
      #3
    \end{Form}

    \par\smallskip
    \switchocg{sol-#1}{\fbox{Mostrar / ocultar solución}}%
    \hfill{\footnotesize #6}\par\smallskip

    \begin{ocg}{Solución}{sol-#1}{0}
      \noindent #4
    \end{ocg}
  \par\medskip
}

%% ------------------- Listas----------------------%%
    % --- Lista alfabética ---
    \newenvironment{la}
      {\begin{list}{\alph{alphaitem})}{%
          \usecounter{alphaitem}%
          \setlength{\leftmargin}{1cm}%
          \setlength{\parskip}{3pt}% 
          \setlength{\itemsep}{0pt}% ← espacio entre ítems
          \setlength{\parsep}{0pt}%  ← párrafos dentro de un ítem
      }}
      {\end{list}}
      
    % --- Lista numérica ---
    \newenvironment{lnum}
      {\begin{list}{\arabic{numitem}.}{%
          \usecounter{numitem}%
          \setlength{\leftmargin}{1cm}%
          \setlength{\parskip}{3pt}% 
          \setlength{\itemsep}{0pt}% ← espacio entre ítems
          \setlength{\parsep}{0pt}%  ← párrafos dentro de un ítem
      }}
      {\end{list}}
      
% ------------------- Imágenes----------------------%
    \graphicspath{{Img/}}% Rutas por defecto 
    % --- Imagen adaptativa ---
    % Registros de dimensión definidos una sola vez
    \newdimen\imgcap
    \newdimen\imgmin
    \newdimen\imgmax
    \newdimen\imgremain
    \newdimen\imgh
    \newdimen\imgneed
    
    \newcommand{\imagenfit}[3]{%
      \addvspace{10pt}% separación antes
      \par\begingroup
        % Parámetros de composición (asignaciones, no \newdimen)
        \imgcap=1\baselineskip            % alto estimado de título+fuente
        \imgmin=0.15\textheight
        \imgmax=0.15\textheight
        \imgremain=\pagegoal
        \ifdim\pagegoal=\maxdimen \imgremain=0pt\fi
        \advance\imgremain by -\pagetotal
        \advance\imgremain by -\imgcap
        % Decisión de altura destino
        \ifdim\imgremain<\imgmin
          \newpage
          \imgh=0.20\textheight
        \else
          \imgh=\imgremain
          \ifdim\imgh>\imgmax \imgh=\imgmax \fi
        \fi
        % Reservar espacio para evitar cortes del bloque completo
        \imgneed=\imgh \advance\imgneed by \imgcap
        \Needspace{\imgneed}
        % Cuerpo
        \noindent\begin{minipage}{\textwidth}
          \centering
          \captionof{figure}{\textemdash\ #2}\par\vspace{0.3em}% título arriba
          \includegraphics[width=\linewidth,height=\imgh,keepaspectratio]{\detokenize{#1}}\par
          \vspace{0.3em}\small\textit{Fuente: }#3% fuente abajo
        \end{minipage}%
      \endgroup
      \par\addvspace{10pt}%
    }

%------------ Bloque de RECUADRO ESPECIAL -------------%
    \newenvironment{specialblock}[1]{%
      \begin{quote} % crea sangría a izquierda y derecha
      \small % texto más pequeño
      \noindent\textbf{#1}\\[-0.3em] % título en negrita
      \rule{\linewidth}{0.4pt}\\[0.6em]}{
      \end{quote}}
      
%-------------------------- Portada -----------------------%
\newcommand{\oldtitle}[1]{\def\OldTitle{#1}}
\newcommand{\changerationale}[1]{\def\ChangeWhy{#1}}

\newcommand{\changebox}{%
\begin{tcolorbox}[title={Historial de cambios del tema}]
\textbf{Título anterior:} \OldTitle\par
\textbf{Motivo del cambio:} \ChangeWhy
\end{tcolorbox}
}

%-------------------------- Author Font -----------------------%
    \newcommand{\authorfont}[1]{{\fontfamily{EBGaramond-TLF}\selectfont #1}}

%-------------------------- Anexos -----------------------%
    \makeatletter
    \newcounter{anexo}
    \renewcommand{\theanexo}{\Roman{anexo}}
    \newcommand{\anexo}[1]{%
      \clearpage
      \phantomsection
      \refstepcounter{anexo}%
      \noindent\textbf{Anexo \theanexo.- }#1\par\medskip
      \addcontentsline{stc}{section}{Anexo \theanexo.- #1}%
    }
    \makeatother

    %-------------------------- Lista de figuras (personal) -----------------------%
\makeatletter
\newcommand{\MyListOfFigures}{%
  \clearpage
  \phantomsection
  {\centering\bfseries\Large Índice de figuras\par}%
  \medskip
  % Entrada en nuestro índice de contenido (stc)
  \addcontentsline{stc}{section}{Índice de figuras}%
  % Recuperación del listado (sin cabecera propia)
  \begingroup
    \parindent=0pt\parskip=2pt
    \@starttoc{lof}%
  \endgroup
}
\makeatother

%-------------------------- Bibliografía -----------------------%
\makeatletter
% Contador para las referencias
\newcounter{myref}
% Cabecera de la bibliografía, entra en el índice 'stc'
\newcommand{\MyBibliography}{%
  \clearpage
  \phantomsection
  {\centering\bfseries\Large Bibliografía y referencias\par}%
  \medskip
  \addcontentsline{stc}{section}{Bibliografía y referencias}%
  \setcounter{myref}{0}%
}
\newcommand{\MyBibItem}[4]{%
  \refstepcounter{myref}%
  \noindent[\themyref]\ #1.\ \emph{#2}. #3. #4\par
}
\makeatother

%--------- Establecer cuatro niveles de índice --------%
    \setcounter{tocdepth}{4}   
    \setcounter{secnumdepth}{4}% numera hasta \paragraph

%%%%%%%%%%%%%%%%%%%%%%%%%%%%%%%%

\makeatletter

    %--------- Interlineado Global --------%
    \edef\GlobalBaseLineStretch{\baselinestretch}
    
    %--------- Espaciado de seccinones --------%
    \renewcommand\section{\@startsection{section}
      {1}{\z@}
      {2.5ex \@plus 1ex \@minus .2ex} % espacio antes
      {1.5ex \@plus .2ex}             % espacio después
      {\normalfont\Large\bfseries}    % formato del título
    }
    \renewcommand\subsection{\@startsection{subsection}{2}{\z@}
      {3ex \@plus 0.8ex \@minus .2ex}
      {1.5ex \@plus .2ex}
      {\normalfont\large\bfseries}
    }
    \renewcommand\subsubsection{\@startsection{subsubsection}{3}{\z@}
      {3ex \@plus 0.5ex \@minus .2ex}
      {1ex \@plus .2ex}
      {\normalfont\normalsize\bfseries}
    }
    \renewcommand\paragraph{\@startsection{paragraph}{4}{\z@}%
    {-2ex \@plus -0.5ex \@minus -.2ex}% espacio antes
    {0.8ex \@plus .2ex}% espacio después
    {\normalfont\normalsize\itshape}}% ← cursiva en lugar de negrita

    %--------- Ecuaciones --------%   
    % Variables globales para almacenar títulos y notas
    \gdef\leq@current@section{}%
    \gdef\leq@current@subsection{}%
    \gdef\leq@last@printed@section{}%
    \gdef\leq@last@printed@subsection{}%
    
    % Comando para añadir nota (invisible en el cuerpo)
    \newcommand{\eqnote}[1]{%
      \addcontentsline{leq}{leqnote}{#1}%
    }
    
    % Comando para ecuaciones
    \newcommand{\eqblock}[2]{%
      % Verificar si necesitamos imprimir el título de sección
      \ifx\leq@current@section\leq@last@printed@section\else
        \addcontentsline{leq}{leqsection}{\leq@current@section}%
        \global\let\leq@last@printed@section\leq@current@section
        \gdef\leq@last@printed@subsection{}% Reset subsección al cambiar sección
      \fi
      % Verificar si necesitamos imprimir el título de subsección
      \ifx\leq@current@subsection\leq@last@printed@subsection\else
        \ifx\leq@current@subsection\@empty\else
          \addcontentsline{leq}{leqsubsection}{\leq@current@subsection}%
          \global\let\leq@last@printed@subsection\leq@current@subsection
        \fi
      \fi
      % Ecuación
      \refstepcounter{eqnum}%
      \begin{equation*}
        #1\tag{E.\theeqnum}
      \end{equation*}%
      \addcontentsline{leq}{leqt}{[E.\theeqnum] -- #2}%
      \addcontentsline{leq}{leqm}{#1}%
    }
    
    %%% ----- Índice de ecuaciones ----------%%%
    % Tamaño para []
    \newcommand{\leqIndexFont}{\normalsize}
    \newcommand{\leqTitleFont}{\tiny}
    \newcommand{\leqNoteFont}{\tiny}
    % Cabecera de sección 
    \newcommand*\l@leqsection[2]{%
      \addvspace{25pt}% Mayor separación antes de nueva sección
      \noindent\leqIndexFont
      \makebox[\linewidth]{\rule{.38\linewidth}{0.4pt}\ \ \textbf{  \MakeUppercase{#1}}\ \ \rule{.38\linewidth}{0.4pt}}%
      \par\vspace{-10pt}
    }
    % Cabecera de subsección 
    \newcommand*\l@leqsubsection[2]{%
      \par\addvspace{15pt}%
      \noindent\leqIndexFont #1
      \par\vspace{-7pt}
    }
    % Título de ecuación 
    \newcommand*\l@leqt[2]{%
    \par\addvspace{10pt}%
      \par\noindent\leqTitleFont #1\par
    }
    % Miniatura de ecuación
    \newcommand*\l@leqm[2]{%
      \vspace{0.3\baselineskip}%
      \par\scriptsize\noindent\hspace*{1.5em}\(#1\)\par%
    }
    % Nota de ecuación 
    \newcommand*\l@leqnote[2]{%
      \vspace{0.2\baselineskip}%
      \par\noindent\leqNoteFont\hspace*{1.5em}\textbf{Nota.--} #1\par\vspace{0.5\baselineskip}%
    }
    % Generación del índice
    \newcommand{\mylistofequations}{%
      \clearpage
      \phantomsection
      % Título visual de la lista (ajústalo a tu gusto):
      {\centering\bfseries\Large Índice de ecuaciones\par}
      % Entrada en nuestro índice 'stc':
      \addcontentsline{stc}{section}{Índice de ecuaciones}%
      % Llamada a tu lista real (usa el nombre exacto de tu macro):
      \@starttoc{leq}%
    }
    
    %--------- Captura de títulos de sección --------%
    \let\leq@original@section\section
    \let\leq@original@subsection\subsection
    % Flag para saber si estamos en una sección asterisco
    \newif\ifLeqStarSection
    \LeqStarSectionfalse
    
    % Redefinir \section CON CONTROL DE ASTERISCO
    \renewcommand{\section}{%
      \@ifstar{\leq@section@star}{\@dblarg\leq@section@nostar}%
    }
    \def\leq@section@star#1{%
      \global\LeqStarSectiontrue
      \gdef\leq@current@section{#1}%
      \gdef\leq@current@subsection{}%
      \leq@original@section*{#1}%
      \global\LeqStarSectionfalse
    }
    \def\leq@section@nostar[#1]#2{%
      \global\LeqStarSectionfalse
      \gdef\leq@current@section{#2}%
      \gdef\leq@current@subsection{}%
      \leq@original@section[#1]{#2}%
    }
    % Redefinir \subsection
    \renewcommand{\subsection}{%
      \@ifstar{\leq@subsection@star}{\@dblarg\leq@subsection@nostar}%
    }
    \def\leq@subsection@star#1{%
      \global\LeqStarSectiontrue
      \gdef\leq@current@subsection{#1}%
      \leq@original@subsection*{#1}%
      \global\LeqStarSectionfalse
    }
    \def\leq@subsection@nostar[#1]#2{%
      \global\LeqStarSectionfalse
      \gdef\leq@current@subsection{#2}%
      \leq@original@subsection[#1]{#2}%
    }

    % ----- Restauración de valores Globales de estilo ----------%
    % Flags
    \newif\ifInFootnote
    \InFootnotefalse
    \pretocmd{\@footnotetext}{\InFootnotetrue}{}{}
    \apptocmd{\@footnotetext}{\InFootnotefalse}{}{}
    % Profundidad de listas (soporta anidadas)
    \newcount\MyListDepth \MyListDepth=0
    \AtBeginEnvironment{la}{\advance\MyListDepth by 1}
    \AfterEndEnvironment{la}{\advance\MyListDepth by -1}
    \AtBeginEnvironment{lnum}{\advance\MyListDepth by 1}
    \AfterEndEnvironment{lnum}{\advance\MyListDepth by -1}
    \AtBeginEnvironment{itemize}{\advance\MyListDepth by 1}
    \AfterEndEnvironment{itemize}{\advance\MyListDepth by -1}
    \AtBeginEnvironment{enumerate}{\advance\MyListDepth by 1}
    \AfterEndEnvironment{enumerate}{\advance\MyListDepth by -1}
    % Restauración diferida: hazlo al INICIO del próximo párrafo real
    \newif\if@pendingRestore
    \@pendingRestorefalse
    \newcommand{\RequestRestore}{\global\@pendingRestoretrue}
    \everypar{%
      \if@pendingRestore
        \global\@pendingRestorefalse
        \ifInFootnote\else
          \ifnum\MyListDepth=0
            \setlength{\parskip}{\GlobalParskip}%
            \setstretch{\GlobalBaseLineStretch}%
          \fi
        \fi
      \fi
    }
    % Al final de bloques:
    \AfterEndEnvironment{equation}{\RequestRestore}
    \AfterEndEnvironment{equation*}{\RequestRestore}
    \AfterEndEnvironment{la}{\RequestRestore}
    \AfterEndEnvironment{lnum}{\RequestRestore}
    % Al final de macros:
    \apptocmd{\eqblock}{\RequestRestore}{}{}
    \apptocmd{\imagenfit}{\RequestRestore}{}{}
    
\makeatother

% ===================== ToC propio con 4 niveles (único bloque) =====================
\makeatletter

% --- Escritores: añadimos una línea al .stc en cada cabecera numerada ---
% Guardamos las versiones actuales para llamar al comportamiento normal
\let\MyTOC@orig@section\section
\let\MyTOC@orig@subsection\subsection
\let\MyTOC@orig@subsubsection\subsubsection
\let\MyTOC@orig@paragraph\paragraph

% SECTION
\renewcommand{\section}{%
  \@ifstar{\MyTOC@section@star}{\@dblarg\MyTOC@section@nostar}%
}
\def\MyTOC@section@star#1{%
  \MyTOC@orig@section*{#1}% sin entrada al .stc
}
\def\MyTOC@section@nostar[#1]#2{%
  \MyTOC@orig@section[{#1}]{#2}% comportamiento normal (escribe al .toc)
  % Escribimos también nuestra línea al .stc (texto con \numberline)
  \addcontentsline{stc}{section}{\protect\numberline{\thesection}#2}%
}

% SUBSECTION
\renewcommand{\subsection}{%
  \@ifstar{\MyTOC@subsection@star}{\@dblarg\MyTOC@subsection@nostar}%
}
\def\MyTOC@subsection@star#1{%
  \MyTOC@orig@subsection*{#1}%
}
\def\MyTOC@subsection@nostar[#1]#2{%
  \MyTOC@orig@subsection[{#1}]{#2}%
  \addcontentsline{stc}{subsection}{\protect\numberline{\thesubsection}#2}%
}

% SUBSUBSECTION
\renewcommand{\subsubsection}{%
  \@ifstar{\MyTOC@subsubsection@star}{\@dblarg\MyTOC@subsubsection@nostar}%
}
\def\MyTOC@subsubsection@star#1{%
  \MyTOC@orig@subsubsection*{#1}%
}
\def\MyTOC@subsubsection@nostar[#1]#2{%
  \MyTOC@orig@subsubsection[{#1}]{#2}%
  \addcontentsline{stc}{subsubsection}{\protect\numberline{\thesubsubsection}#2}%
}

% PARAGRAPH
\renewcommand{\paragraph}{%
  \@ifstar{\MyTOC@paragraph@star}{\@dblarg\MyTOC@paragraph@nostar}%
}
\def\MyTOC@paragraph@star#1{%
  \MyTOC@orig@paragraph*{#1}%
}
\def\MyTOC@paragraph@nostar[#1]#2{%
  \MyTOC@orig@paragraph[{#1}]{#2}%
  % Si no numeras los \paragraph, cambia \theparagraph por texto plano sin \numberline
  \addcontentsline{stc}{paragraph}{\protect\numberline{\theparagraph}#2}%
}

% --- Impresor del índice propio (lee el .stc) ---
\newcommand{\MyTOC}{%
  \par\bigskip
  {\centering\bfseries\Large Índice de contenido\par}%
  \medskip
  \begingroup
    \parindent=0pt\parskip=2pt
    \@starttoc{stc}%
  \endgroup
}

\makeatother
% ===================== fin ToC propio =====================


%%%%%%%%%%%%%%


% --- Datos del tema ---
\title{Unión Europea (VII): Relaciones económicas exteriores de la UE\\
\large La política de cooperación al desarrollo de la Unión Europea}
\author{Marcelo Ortega}
\date{Julio 2026}
\newcommand{\TemaCode}{3B42}

\oldtitle{
    B.44. Las relaciones económicas exteriores de la Unión Europea. La política de cooperación al desarrollo de la Unión Europea
}
\changerationale{
    Sin cambios.
}

\usepackage{soul}\sethlcolor{yellow}% resaltado amarillo: \hl{texto} (Panel TCEE, ⌃H)
\begin{document}


\maketitle
\changebox

\newpage

% --- Página de resumen y modificaciones ---
\puntosclave{ %% Escribir puntos clave Apuntes CECO
     Para abordar satisfactoriamente este título del temario, se consideran fundamentales dos decisiones: (i) la delimitación del contenido, (ii) la aproximación a seguir. Sobre el contenido: definiremos el ámbito de estudio negativamente respecto al de otros temas de la UE. En primer lugar, se trata de un tema acerca de las relaciones de la UE frente al resto del mundo; por ello, resulta importante no incluir cuestiones relacionadas con el mercado interior (en otras palabras, se ha de tratar el movimiento de factores inter-UE y otras economías, pero no intra-UE, a diferencia de lo que se espera en el tema 3839). En segundo lugar, dada la existencia del tema 3841 para estudiar con detalle las relaciones comerciales de la UE, en éste no se incluirá lo relacionado con comercio internacional de bienes y servicios.

     Sobre la aproximación: se opta por una combinación entre el fonnalismo jurídico e institucional de los demás títulos de UE en el temario y un enfoque de Economía política. Es decir, han de entenderse las relaciones económicas y la política económica exterior como instrumentos geopolíticos que la Unión aplica para defender sus intereses internacionalmente, y conseguir así que prevalezcan su visión y valores. Es lo que en Economía se denomina frecuentemente como \textbf{diplomacia económica}.

     [De los autores del Tema de CECO]. Consideramos adecuado estudiar en torno a tres ejes diferenciados. Por un lado, cómo la UE trata de consagrar un área de prevalencia en la región (expandiendo sus fronteras e influyendo sobre sus vecinos), a la vez que defiende sus posturas en foros multilaterales y se garantiza un grado de autonomía suficiente a nivel internacional. Por otro lado, cómo los intercambios económicos nocomerciales (i.e., movimiento de factores) detenninan las relaciones de la UE con el resto del mundo en general, y cómo son esas relaciones con algunos actores geopolíticos relevantes en particular. Finalmente, la política de cooperación al desarrollo, parcela en la que el modelo europeo prueba su carácter generoso y humanista frente a otros, pero que también contiene importantes aspectos estratégicos.
    }
\vspace{1cm}

\modificaciones{
    
    }

\newpage
\MyTOC
\newpage

\mylistofequations
\clearpage

\MyListOfFigures
\clearpage


\pagenumbering{arabic}
\setcounter{page}{1}


\section{Introducción}



\subsection{Contextualización}


\subsubsection{Relevancia}




\subsubsection{Problemática}



\subsection{Estructura de la exposición}




\clearpage
\section{Acción Exterior de la Unión: Marco jurídico-conceptual}
\puntosclave{
    La perspectiva de la adhesión a la UE tiene un poderoso efecto transformador en los aspectos democrático, político, económico y social para los países que se embarcan en este proceso. Además de la evolución de estos países, la adhesión de nuevos miembros también genera nuevas oportunidades y retos para la UE.
    
    La Política Europea de Vecindad basada en el compromiso mutuo con unos valores comunes entre la UE y sus países vecinos busca fortalecer la prosperidad, la estabilidad yla seguridad en sus fronteras.
    
    La UE no sólo ejerce sus relaciones económicas de forma bilateral y regional, sino también mediante su participación en las organizaciones internacionales y foros multilaterales donde defiende y promueve sus intereses en temas climáticos, financieros, etc. Sin embargo, la falta de una posición común entre los EEMM en algunas materias y la imposibilidad de participar directamente (teniendo que hacerlo a través de los EEMM) en algunas de las instituciones limita la efectividad de esta actividad.
}


El proceso de construcción de la Unión Europea cuenta con una dimensión exterior que se muestra como \textbf{uno de sus componentes más singulares y complejos}.\footnote{%
    Como advirtió tempranamente STEIN, las tensiones existentes en el marco de las relaciones entre los Estados y las CCEE y los conflictos de separación de poderes entre las distintas Instituciones \textit{adquieren un matiz diferente cuando se trata de las relaciones con el mundo exterior} (1990: 127). Certera apreciación: todo lo que tiene que ver con el mundo de la acción exterior en la UE parece obedecer a una lógica distinta.
} En efecto, \textbf{situada en el núcleo más sensible de la soberanía de los Estados}, la política exterior ofrece mayores resistencias como objeto del fenómeno de integración. 

Pese a ello, la construcción europea no podía carecer de una dimensión exterior, esencialmente, por tres razones: la imposibilidad de adoptar una política de autarquía en un contexto internacional de interdependencia, los vínculos coloniales que mantenían algunos de sus Estados miembros y, naturalmente, la repercusión misma de la creación de este nuevo mercado comunitario en las relaciones comerciales internacionales (BOURRINET y TORRELLI, 1992: 4). Por tanto, el proceso de integración supuso la atribución de competencias en materia de relaciones económicas exteriores que implicaba, a su vez, la separación entre los aspectos económicos y el resto de los ámbitos de política exterior. 

Esa distinción, tan lógica en el plano teórico como artificial en la práctica, lleva a que hablemos de \textbf{dos formas de acción exterior}: (i) la Política Exterior y de Seguridad Común (\textbf{PESC}), mantenida bajo una lógica intergubernamental y articulada mediante un sistema más cercano a la concertación donde la presencia de lo estatal es más condicionante; y, (ii) las relaciones exteriores fruto de la Acción Exterior competencia de la Unión Europea.\footnote{%
    \textbf{NOTA.-} Hasta ahora esta separación entre la acción exterior de la Unión Europea y la de PESC se asentaba en la radical apreciación de un núcleo de relaciones exteriores fundadas en la personalidad jurídica de las CCEE y en la carencia de PJI de la UE que sustentaba la acción exterior de la PESC y de la CPJP. No obstante, con la inclusión del art. 37 TUE que permite a la Unión celebrar acuerdos con terceros tanto en política exterior y de seguridad como en el pilar de la cooperación policial y judicial penal, se empieza a plantear la necesidad de revisar aquella radical apreciación.
} 

Esta lógica sigue manteniéndose hoy en día, pese a la cada vez más intrincada relación entre ambas y la problemática asociada: el TUE sigue ofreciendo un capítulo separado (Cap. 2, Título V) que establece las disposiciones específicas para la PESC, mientras que la acción exterior de la Unión que denominamos relaciones exteriores queda encuadrada en el TFUE (Quinta Parte).


\subsection{Elementos Comunes: ¿Qué es la Unión Europea en el ámbito internacional?}
No obstante, a pesar de esta dicotomía, el proyecto de integración europeo, como proyecto que excede del ámbito exclusivamente económico y atendiendo a los objetivos que persigue, se ha esforzado por establecer una serie de elementos o directrices comunes que deben vertebrar la acción exterior de la Unión, en tanto que ente independiente, pero también de sus Estados Miembros, en tanto que sujetos ligados a los objetivos de la Unión.

\subsubsection{Principios de Acción Exterior: TUE (Tít. V, Cap. 1º)}
El primero de éstos, y más importante, es el conjunto de disposiciones generales previstos en el TUE (Título V, Capítulo 1º, arts. 21 y 22), que, estableciendo los \textbf{principios orientadores de la acción exterior de la UE} e \textbf{indentificando los objetivos de las políticas comunes y acciones de la UE en materia exterior}, introducen cierta coherencia y unidad en todas las acciones exteriores (institucional, procedimental y material).

En particular, el art. $21$ TUE dicta los principios que orientan la acción exterior de la UE, que son la democracia, el Estado de Derecho y el respeto de los derechos humanos, de conformidad con los principios del Derecho Internacional, e identifica como objetivos de la acción exterior la defensa de sus valores, intereses y seguridad, junto con el mantenimiento de la paz internacional o el apoyo a los países en desarrollo y el fomento del comercio internacional y el medio ambiente. 

Por último, establece que la Unión velará por la coherencia entre los distintos ámbitos de la acción exterior y entre ésta y las demás políticas, responsabilizando de esa coherencia al Consejo y la Comisión en lo que respecta a la Acción Exterior, y en lo relativo a la PESC responsabilizando al Consejo, teniendo un papel más marginal la Comisión. Y, en todo caso, asistidos por el elemento institucional que es la figura del Alto Representante para Asuntos Exteriores y Política de Seguridad.\footnote{%
    Es interesante resaltar com el art. 22 TUE, en ese ímpetu por elevar al ámbito comunitario la acción exterior en su conjunto, otorga un papel al Alto Representante y a la Comisión en todo lo relativo a la Acción Exterior de la UE, aunque sea residual en el ámbito de la PESC.
}



\subsubsection{¿Unión Europea como sujeto internacional? Personalidad Jurídica Internacional}
El segundo de estos elementos comunes que vertebran la acción exterior de la Unión es su condición como sujeto de derecho internacional, que queda perfectamente reconocida en el art. $47$ TUE: \textit{la Unión tiene personalidad jurídica}. 

Ahora bien, ¿de dónde emerge y qué clase de personalidad jurídica tiene? La doctrina es unánime en este respecto: la Unión Europea, como todas las OI, goza de personalidad jurídica derivada, es decir funcional, condicionada por la voluntad de sus Estados miembros y limitada en cuanto a su alcance y contenido.\footnote{%
    Las normas de Derecho Internacional público relativas a la subjetividad jurídico-internacional establecen una distinción nítida y rotunda entre los Estados como sujetos originarios y las organizaciones internacionales que gozan de una personalidad jurídica derivada, condicionada por la voluntad de sus Estados miembros, que encuentra su expresión jurídica en el consentimiento, y determinada también por la función que se atribuye a la organización. La PJI que ostenta la UE es, por tanto, \textit{igual}, de iure, al resto de las organizaciones.
} Pero aquí empiezan a surgir problemas. 

Las funciones asignadas a la UE exceden con mucho del conjunto de atribuciones generalmente conferidas a las OI clásicas, que cumplen una función de coordinación de intereses y operan sobre una limitación de los poderes estatales. En cambio, las organizaciones de integración que tienen atribuida una función de unificación se alimentan de poderes de los Estados, de manera que la limitación es sustituida por una \textit{transferencia} de poderes. Por esta razón, la UE sigue planteando en el ámbito exterior una serie de problemáticas propias: su \textbf{singularidad} misma como modelo de organización, las \textbf{tensiones} existentes entre la UE y sus Estados miembros (que se agudizan en el ámbito exterior), el \textbf{carácter evolutivo} del proceso de integración, la \textbf{indefinición} de su término final y los condicionantes y límites que impone el Derecho Internacional público como ordenamiento que regula la sociedad internacional.\footnote{%
    Como prueba de estos condicionantes detacan los problemas que genera la delimitación de las competencias exteriores de la Unión y de los Estados miembros, la participación de la UE en tratados, organizaciones y conferencias internacionales o la carencia del \textit{locus standi} en el marco de los procedimientos ante el TIJ que obliga a acudir a otros medios de solución de controversias abiertos a las OI.
} Además, este artículo, que se interpretaba con una cierta tendencia restrictiva, fue extendida progresivamente, en virtud de otras normas (originarias), que concretan el alcance y contenido de esa personalidad y, también, por la interpretación del TJUE.\footnote{%
    De hecho, la jurisprudencia del TJUE ha contribuido considerablemente a la precisión del alcance y contenido de la capacidad internacional de la hoy UE. Una de las más relevantes, y primeras en la materia fue su sentencia del asunto AETR, donde afirmaba con rotundidad que la hoy Unión \textit{goza de la capacidad de establecer vínculos contractuales con Estados terceros en toda la extensión del ámbito de los objetivos definidos en la primera parte del Tratado}.
}

Así, la originalidad de la UE como modelo de organización, a medio camino entre las OI clásicas y las formas estatales de organización, se hace particularmente evidente en el ámbito exterior: aunque su condición de sujeto de Derecho Internacional está fuera de duda, el contenido y alcance de su personalidad jurídica y sus distintas manifestaciones son una expresión privilegiada de la singularidad de la UE como modelo de organización y son, también, el origen de buena parte de los problemas que plantea \textit{ad internum} y \textit{ad externum} su Acción Exterior. 

\paragraph{Manifestación de su PJI: ¿Qué puede hacer la Unión Europea?}
Habida cuenta de lo expuesto, y sin perder la generalidad propia de esta apartado (puesto que estamos hablando de la acción exterior en su conjunto), debemos ver los principales atributos que conforman la PJI de la UE:\footnote{%
    \textbf{NOTA.-} La doctrina entiende que dos de estas capacidad se corresponden con una única naturales: las relaciones de la UE con otras sujetos de derecho internacional. Estas relaciones pueden articularse, básicamente, a través de tres cauces distintos: la representación mediante el ejercicio del derecho de legación activo y pasivo, la cooperación o colaboración administrativa entre organizaciones internacionales (ex art. 220 TFUE)
}
\begin{la}
    \item \textbf{Capacidad para conluir Tratados: Poder de compromiso}. El poder de concluir tratados internacionales es uno de los atributos que caracteriza la personalidad jurídica de la UE y que sirve para singularizarla respecto del resto de las OI, que están lejos de desarrollar, cualitativa y cuantitativamente, una actividad convencional como la de la Unión;
    \item \textbf{Capacidad de respresentarse: Derecho de legación (activo/pasivo)}. El derecho de legación activo consiste en la facultad de enviar representantes ante terceros Estados y organizaciones internacionales, mientras que el derecho de legación pasivo es el derecho de recibir agentes diplomáticos o representantes de otros Estados y OI. Con la Decisión de julio de 2010, el Consejo estableció la organización y el funcionamiento del Servicio Europeo de Acción Exterior (SEAE), con la intención de integrar los servicios de la Comisión y del Consejo implicados en la acción exterior, dando lugar a un servicio sui generis, funcionalmente autónomo, cercano así a una institución, dependiente del Alto Representante, cuya función consiste en representar a la UE ante terceros Estados y OOII;\footnote{%
        Recogida en el art. 27 TUE, se prevé un Servicio Europeo de Acción Exterior, aunque no existe ninguna disposición explícita relativa a ninguno de los dos derechos. A pesar de ello, la UE han venido ejerciendo ampliamente ambos derechos con una \textbf{fundamentación distinta en cada caso}. Por un lado, el \textbf{derecho de legación pasivo} ha sido fruto del trabajo del Consejo y la Comisión en definir las competencias y los procedimientos en relación con la acreditación de misiones y representaciones de terceros ante la UE. No exento de problemas (\textit{Compromiso de Luxemburgo}),  el sistema se articula exigiendo el consentimiento unánime de los Estados miembros para el establecimiento de la misión y la presentación de cartas credenciales, separadamente, en el mismo día y sin ceremonial, al Presidente del Consejo y al Presidente de la Comisión. Por su parte, el \textbf{derecho de legación activo} encuentra su fundamento en la práctica de la Comisión consistente en el establecimiento de oficinas o delegaciones exteriores sobre la base del poder de organización interna de esta Institución y, naturalmente, de la aceptación de los terceros Estados y OOII. Como tampoco tardaron en aparacer los problemas, la cuestión se resolvió aceptando una doble representación, que correspondía a la Comisión y al Consejo que nos trae a la actualidad: la función de representación es ejercida por las delegaciones de la Unión (art. 221 TFUE) y por las misiones diplomáticas de los Estados miembros.
    }
    \item \textbf{Capacidad para participar: la UE y las Organizaciones y Conferencias Internacionales}. La participación de la UE en organizaciones y conferencias internacionales constituye una manifestación singular del alcance y contenido de su personalidad jurídica internacional. Como es normal, esta participación está sujeta a tres condicionantes: (i) \textbf{alcance y contenido}, como el Derecho originario de la Unión carece de una disposición expresa sobre la participación de la UE, la mayoría de la doctrina entiende que los artículos que confieren competencia expresa a la UE para celebrar tratados internacionales o cooperar con terceros Estados u organizaciones internacionales constituyen base jurídica suficiente para que la UE se convierta en miembro de otras organizaciones internacionales; (ii) \textbf{objeto y finalidad} de la organización o conferencia en cuestión, pues su participación depende a menudo de los tratados constitutivos de estas OOII, que dificultan o impiden que la UE se convierta en miembro de ellas, al contemplar sólo la posibilidad de que los Estados accedan a la condición de miembros;\footnote{%
        No obstante, la UE ha sido pionera en la ruptura de la intergubernamentalidad de las organizaciones internacionales y ha conseguido que los tratados constitutivos de algunas de ellas le permitan ser miembro, como es el caso de la FAO (PAC) o de la OMC (PCC). Teniendo en cuenta los dos primeros condicionantes, la participación de la UE en las mismas presenta las siguientes variantes: (i) \textbf{UE con estatuto de observador}, que abarca una amplia casuística; (ii) \textbf{UE con membresía exclusiva}, normalmente en aquellos organismos, de carácter técnico, que gestionan materias en las que la UE tiene competencias exclusivas (muy pocas); y, (iii) \textbf{UE con membresía conjunta}, participando conjuntamente con sus Estados miembros, normalmente por la proliferación de los llamados \textit{acuerdos mixtos} y siempre en organizaciones que actúan sobre materias de competencia compartida de la UE.
    
    En este último caso, como las competencias exclusivas de la UE son excepcionales y la regla general es la existencia de competencias compartidas o concurrentes, parece lógico que se afirme progresivamente la participación conjunta de la UE y de sus Estados miembros en las organizaciones internacionales. No obstante, aparece una complicación adicional: determinar el alcance de sus competencias respectivas y proceder a su necesaria coordinación en el marco de la organización o de la conferencia internacional. Como es normal, no está exenta de cierta conflictividad y deriva hacia su resolución judicial (Comisión c. Consejo, 1995 TJUE, por ejemplo).
    } y, (iii) \textbf{régimen de participación} previsto en la organización o conferencia, que se ha venido articuladondo mediante delegación bicéfala, delegación única y la representación por la Comisión.
\end{la}

Hasta aquí el conjunto de elementos comunes que tienen ambas ramas o vertientes de la acción exterior de la Unión Europea.

\subsection{Acción Exterior de la Unión: Naturaleza y alcance}
Empezamos por presentar la naturaleza, alcance, objetivos y competencias de la Unión en virtud de los dispuesto en la Quinta Parte del TFUE, que llamaremos: la Acción Exterior de la Unión Europea. 

\subsubsection{Tratado de Funcionamiento de la Unión Europea (2007): Articulación jurídica actual}
La Quinta Parte del TFUE agrupa una serie de disposiciones hasta ahora dispersas y que regulan los principales ámbitos materiales de carácter puramente exterior de la Unión (como la política comercial o la cooperación al desarrollo) junto con las competencias instrumentales (como la celebración de tratados, la relación con otras organizaciones internacionales o la adopción de medidas restrictivas). Se trata, como veremos, de un conjunto de competencias rebelde frente a cualquier tentativa de sistematización, cuya delimitación se presenta como una operación extraordinariamente compleja y que solo ha sido parcialmente solventada con el Tratado de Lisboa introduciendo una cierta coherencia.\footnote{%
    Fundamentalmente por dos razones. En primer lugar, no se limitan a las expresamente atribuidas en el texto de los Tratados, sino que, junto a ellas, se encuentran las derivadas de la construcción jurisprudencial realizada por el TJUE sobre la base de la teoría de los poderes implícitos y de la utilización del artículo 352 TFUE. En segundo lugar, estos tres cauces de generación de competencias no tienen un carácter excluyente entre sí.
}

\paragraph{Espacio de Libertad, Seguridad y Justicia (ELSJ): \textit{ex} art. 67 TFUE}
Además, y en aras de proporcionar una visión más completa de la Acción Exterior de la Unión Europea, debemos comentar la incidencia que tiene sobre esta el llamado Espacio de Libertad, Seguridad y Justicia (ELSJ), uno de los objetivos de la Unión que mayor atención requieren en la actualidad. Un objetivo, de primera magnitud, que no sólo engloba ámbitos competenciales muy amplios, sino que afecta a sectores de la actividad estatal muy sensibles, política y jurídicamente ya que se proyecta sobre ejes centrales de la vida ciudadana como son el orden penal, la jurisdicción de jueces y tribunales, las labores de las fuerzas de policía o las políticas de asilo o de inmigración.  

La articulación jurídica se encuentra en el TFUE (Tercera Parte, Tít. V), a raíz de la reforma impulsada por el Tratado de Lisboa, que trata de trasladar su enorme trascendencia (aunque su descripción no haya cambiado prácticamente desde el Tratado de Amsterdam). Suelen identificarse tres componentes como los contenidos principales de ese espacio público europeo: (i) la ausencia de controles de personas en las fronteras interiores y el desarrollo de una política común de asilo, inmigración y control de las fronteras;\footnote{%
    \textbf{NOTA.-} Es importante retener que la supresión de los controles en las fronteras interiores no supone el otorgamiento a los ciudadanos europeos del derecho de libre circulación (cuyo marco es el estatuto de ciudadanía y las libertades del mercado interior), pero sí provoca la de los nacionales de terceros Estados, lo que obliga, se dice, a responder con medidas en materia de asilo y flujos de migración.
} (ii) esforzarse por garantizar un nivel elevado de seguridad mediante medidas de prevención y lucha contra la delincuencia, el racismo y la xenofobia, medidas de coordinación y cooperación entre autoridades policiales y judiciales y otras autoridades competentes, así como mediante el reconocimiento mutuo de las resoluciones judiciales en materia penal y, si es necesario, mediante la aproximación de las legislaciones penales; y, (iii) facilitar la tutela judicial, garantizando en especial el principio de reconocimiento mutuo de las resoluciones judiciales y extrajudiciales en materia civil.\footnote{%
    Como vemos, esta descripción general de la estructuración normativa del ELSJ confirma la conocido crítica doctrinal de fondo: una amalgama de objetivos políticos difícilmente reconducibles a unidad y un peso excesivo del vector securitario en detrimento de una casi ausente concepción de la libertad.
}

Como es normal, nos centraremos por su incidencia económica en las Políticas sobre control de fronteras, asilo e inmigración, que el TFUE sistematiza en tres grandes políticas: conjugar la eliminación en las interiores con un control eficaz en el cruce de las exteriores con miras a lograr un progresivo \textbf{sistema integrado de gestión de las fronteras exteriores}. 

En primer lugar, el art. $77$ recoge unas amplias bases jurídicas que posibilitan adoptar medidas para establecer la ausencia total de controles de las personas, sea cual sea su nacionalidad, cuando crucen las fronteras interiores, una política común de visados y otros permisos de residencia de corta duración, las condiciones en las que los nacionales de terceros países podrán circular libremente por la Unión durante un corto período, los controles a los cuales se someterá a las personas que crucen las fronteras exteriores y cualquier otra medida necesaria para el establecimiento progresivo de un sistema integrado de gestión de las fronteras exteriores.\footnote{%
    En este sentido, el Lisboa (2007) en sí no supuso cambio sustancial: en virtud de las normas anteriormente vigentes, ya se había adoptado un denso conjunto normativo sustantivo (con el Código de Fronteras Schengen y el Código Comunitario sobre Visados, a la cabeza), que se acompaña de la infraestructura técnica necesaria para su correcta aplicación, en donde el desarrollo del Sistema de Información Schengen de segunda generación (SIS II) y el Sistema de Información de Visados (VIS) tienen un papel central.

Como sabemos, todo este sistema ha sufrido una presión creciente donde las deficiencias en el control de las fronteras exteriores (incluidos los limitados poderes y recursos de Frontex) se unió a claros incumplimientos de algunos EEMM que, ante la llamada crisis de los refugiados, adoptaron una política de dejar el paso, lo que, a su vez, ha provocado que otros EEMM hayan acudido con largueza a la posibilidad de reinstaurar unilateralmente controles en sus fronteras interiores. La Comisión inició en diciembre de 2015 un proceso para la reconstrucción del espacio Schengen, entre cuyas propuestas se halla el refuerzo y conversión de Frontex en una Agencia de Guardia Europea de Fronteras, sobre la que Consejo y Parlamento parecen haber alcanzado un acuerdo en junio.
} En segundo lugar, con respecto a la reordenación de los temas de asilo, la reforma vino a incluir en el Derecho originario las categorías elaboradas en la normativa derivada previa para la creación de un Sistema Europeo Común de Asilo (SECA), que abarca según fija el art. $78$: un estatuto uniforme de asilo y de protección subsidiaria para los nacionales de terceros países y un sistema común para la protección temporal de las personas desplazadas, en caso de afluencia masiva.\footnote{%
    La preocupación por evitar el \textit{forum shopping} en materia de asilo explica que se prevea igualmente el establecimiento de procedimientos comunes para conceder o retirar el estatuto de asilo o de protección subsidiaria; la fijación de los criterios y mecanismos para determinar qué EM es responsable de examinar una solicitud de asilo o de protección subsidiaria (sistema conocido como Dublín II por sustituir a los Convenios existentes de la era Maastricht) y las normas relativas a las condiciones de su acogida. El SECA se ayuda de la importante base de datos Eurodac donde se incluyen huellas dactilares de solicitantes de asilo y de inmigrantes en situación irregular.
} La revisión de la normativa adoptada durante los primeros años 2000 (Programa de Estocolmo), se hizo imperante pues adolecía de una excesiva laxitud, lo que desembocó en prácticas estatales muy divergentes y en una deficiente implementación. No obstante, La nueva normativa, que suele calificarse como segunda fase del SECA, no parece haber atajado el problema de fondo: la falta de confianza mutua entre los EEMM acerca de un sistema auténticamente común y basado en la solidaridad, como lo prueban los infructuosos intentos de reubicación de refugiados entre EEMM. Finalmente, la Política Común de Inmigración (art. $79$) persigue el tratamiento equitativo de los inmigrantes residentes legales y, particularmente, la lucha contra la inmigración ilegal y la trata de seres humanos.

Además, debemos destacar que el carácter común de estos asuntos explica que estas políticas se rijan por el principio de solidaridad y de reparto equitativo de la responsabilidad entre los EEMM, incluido el aspecto financiero (art. $80$). 

\subsubsection{Régimen competencial: Procedimientos}
Dicho todo esto, ¿cómo se articulan las competencias de la Unión en materia exterior (Acción Exterior) y qué es lo que puede hacer?

La Acción Exterior de la UE encuentra su fundamento en un conjunto de competencias que se ha ido conformando, progresivamente, a través de tres procedimientos sustancialmente distintos: atribución expresa, reconocimiento jurisprudencial y atribución mediante el recurso al artículo 352 TFUE. Habida cuenta de esta situación y en un esfuerzo de ordenación, las competencias exteriores de la Unión se exponen atendiendo a su origen: expresas, implícitas y derivadas del art. 352 TFUE. (Presentaremos únicamente las que van a ser objeto de tratamiento en la segunda sección).

\paragraph{Competencias Expresas}
Las competencias exteriores expresamente atribuidas sobre las que se sustentan las relaciones que trataremos en la exposición son:
\begin{la}
    \item \textbf{La asociación (art. $217$ TFUE)}. La asociación es un modelo de relaciones lo suficientemente amplio para haberse convertido en una figura que dista mucho de ser unitaria o mínimamente uniforme. Siempre supone el establecimiento de relaciones particulares y privilegiadas de la Unión con otros sujetos de derecho internacional. Pero también sirve tanto para preparar una futura adhesión como para establecer marcos específicos de cooperación, que, a su vez, serán diferentes entre sí según se trate de acuerdos con países más desarrollados o con países en vías de desarrollo;
    \item \textbf{Política de Medio Ambiente (art. $191.4$ TFUE)}. La política de medio ambiente constituye una competencia exterior de la Unión sui generis. La determinación del alcance de esta competencia exige tener presente la intervención del principio de subsidiariedad y, por otra parte, la previsión misma de que la cooperación en materia de medio ambiente con otros sujetos se desarrollará por la Unión y por sus Estados miembros en el marco de sus respectivas competencias (art. $191.4$ TFUE);
    \item \textbf{Política Común de Inmigración}. Las disposiciones relativas al ELSJ suponen la comunitarización de ámbitos competenciales que pasan a ser de \textbf{naturaleza compartida}. En aras de la brevedad, aunque el ámbito material cubierto por el ELSJ es vastísimo, lo que a nosotros nos interesa es la política común relativa al control de fronteras, asilo e inmigración, quedando las competencias referidas a: condiciones de entrada y residencia y la expedición de visados y permisos de residencia de larga duración, definición de los derechos de inmigrantes residentes legales en un Estado y las condiciones en que podrán circular y residir en los demás. Además, como es evidente, la puesta en práctica de esta política requiere una estrecha colaboración también con los terceros Estados;\footnote{%
        En este sentido, es cada vez más relevante, política y económicamente, la facultad que ostenta la Unión para la celebración de acuerdos cuyo ámbito material sea la readmisión en sus países de origen o de procedencia, respecto de aquellos que no cumplan o hayan dejado de cumplir las condiciones de entrada, presencia o residencia en el territorio de uno de los EEMM. El recurso al Acuerdo de readmisión firmado con Turquía en diciembre de 2013 y a la aplicación de las excepciones de país seguro y primer país de asilo para resolver la llamada crisis de los refugiados sirios ha sido extraordinariamente contestado política y jurídicamente
    }
    \item \textbf{Cooperación al desarrollo (art. 198 a 204, 208 a 211 y 214 TFUE)}. Bajo esta competencia hemos aglutinado las competencias atribuidas a la Unión en materia de Asociación de los países y territorios de Ultramar, la cooperación al desarrollo, la cooperación económica, financiera y ténica y la ayuda humanitaria. La primer de estas, (arts. 198 a 204), estuvo prevista desde el origen y tenía como finalidad esencial articular las relaciones con países vinculados histórica y económicamente a ciertos Estados europeos (ROY, 1984: 5). Con el tiempo, esas normas se convirtieron en el embrión de un sistema comunitario de cooperación internacional para el desarrollo (ABELLÁN HONRUBIA, 1982: 192), y por esto trataremos todas ellas, conjuntamente, en la última sección.
\end{la}

Se puede apreciar como la enumeración pone de manifiesto la falta de sistemática y de coherencia que preside la articulación en el Tratado de la dimensión exterior del proceso de construcción europea. Además, resulta redundante pero necesario destacar que la mayoría de acciones se materializaran mediante la conclusión de Tratados o Acuerdos, en virtud de la competencia para \textbf{concluir Acuerdos (art. 3.2 TFUE)} y la forma y procedimiento para ejercerlo (art. $218$ TFUE). No hemos mencionado todas, ni si quiera algunas de las más importantes, puesto que son objeto de tratamiento en otras partes del temario.

\paragraph{Competencias implícitas}
La teoría de los poderes implícitos no ha estado ausente en el sistema de la Unión, donde se encuentran tempranamente algunos supuestos de aplicación de esta doctrina en sentencias dictadas por el TJUE. Es, sin embargo, en la década de los años setenta donde esta teoría conoce un desarrollo más significativo y muestra su valor como cauce de extensión de las competencias de la Unión. 

En la actualidad, la doctrina de las competencias implícitas ha sido expresamente consagrada en el art. $216$ TFUE, que ha introducido las consecuencias derivadas de la jurisprudencia del TJUE en relación con la celebración de acuerdos internacionales: \textbf{la competencia exterior emerge cuando la acción de la Unión \textit{sea necesaria para alcanzar, en el contexto de las políticas de la Unión, alguno de los objetivos establecidos en los Tratados}} (\href{https://eur-lex.europa.eu/legal-content/ES/ALL/?uri=CELEX:61976CV0001}{doctrina del Dictamen 1/76}) \textbf{o bien \textit{pueda afectar a normas comunes o alterar el alcance de las mismas}} (\href{https://eur-lex.europa.eu/legal-content/ES/TXT/PDF/?uri=CELEX:61970CJ0022}{doctrina AETR}).\footnote{%
    El origen de la construcción jurisprudencial del TJUE sobre las competencias exteriores implícitas se encuentra en la \href{https://eur-lex.europa.eu/legal-content/ES/TXT/PDF/?uri=CELEX:61970CJ0022}{\textbf{sentencia del asunto AETR}}, donde se afirmaba que las competencias comunitarias podían encontrar su fundamento en el sistema general, en otras disposiciones del Tratado y en actos adoptados en el marco de esas disposiciones. Entendía el TJUE que, una vez que se habían adoptado reglas comunes en el ámbito interno, los Estados miembros no podían asumir compromisos susceptibles de afectar dichas reglas o de alterar su alcance. Un nuevo pronunciamiento, con la \href{https://eur-lex.europa.eu/legal-content/ES/TXT/PDF/?uri=CELEX:61976CJ0003}{\textbf{sentencia del asunto Kramer}}, permite al TJUE avanzar en su interpretación al reconocer la existencia de una competencia comunitaria externa en ausencia de reglas comunes internas. No obstante, es realmente en el \href{https://eur-lex.europa.eu/legal-content/ES/ALL/?uri=CELEX:61976CV0001}{\textbf{Dictamen 1/76, de 26 de abril de 1977}}, donde se encuentra formulada con mayor claridad y amplitud. El TJUE recopilaba su jurisprudencia anterior en la materia para concluir que la competencia exterior derivaba de manera implícita de las disposiciones del Tratado que establecen la competencia interna, siempre que la actuación exterior comunitaria fuese necesaria para la realización de uno de los objetivos de la entonces Comunidad.
}

En el ámbito de las competencias implícitas externa y de carácter exclusivo, fue particularmente importante la sentencia relativa a los conocidos \href{https://eur-lex.europa.eu/ES/legal-content/summary/air-service-agreements-between-eu-and-non-eu-countries.html}{Acuerdos Aereos (2002)}. En esta, el TJUE resolvió acerca de los límites a la capacidad de contraer acuerdos internacionales en materia comercial que tienen los Estados Miembros: \textit{la competencia exclusiva [...] se fundamentaba en que la competencia de la Comunidad para celebrar acuerdos internacionales no sólo es atribuida explícitamente por el Tratado sino que también puede derivarse de otras disposiciones del Tratado y de actos adoptados, en el marco de estas disposiciones, por las Instituciones de la Comunidad}. Con esto, aprovechó para sistematizar su jurisprudencia relativa a las \textbf{condiciones} en que adquiere la Unión una \textbf{competencia externa exclusiva como consecuencia del ejercicio de su competencia interna}, cuando: (i) las obligaciones internacionales están comprendidas dentro del ámbito de aplicación de las normas comunes o dentro de un ámbito ya cubierto en gran medida por tales normas;\footnote{%
    Más recientemente, el TJUE ha señalado en el Dictamen 1/03 que no es necesario que exista una concordancia completa entre el ámbito cubierto por el acuerdo y el de la normativa de la UE ya existente. A estos efectos, el análisis ha de tener en cuenta, no sólo el alcance, sino también la naturaleza y el contenido de las normas comunes, y no sólo el estado actual, sino también las perspectivas de evolución del Derecho de la Unión.
} (ii) la Unión ha incluido en sus actos legislativos internos cláusulas relativas al trato que ha de otorgarse a los nacionales de países terceros o ha conferido expresamente a sus Instituciones competencia para negociar con los países terceros; y, (iii) la Unión haya llevado a cabo una armonización completa en un ámbito determinado, ya que las normas comunes así adoptadas podrían verse afectadas, en el sentido de la sentencia AETR, si los Estados miembros conservaran libertad de negociación con los países terceros. 

Naturalmente, también se ha pronunciada sucesivas veces sobre los límites competenciales en estas materias, que fueron perfectamente manifestadas en el Dictamen 1/94 sobre la OMC (GATS) y, en el caso del Acuerdo sobre Aspectos de los Derechos de Propiedad Intelectual relacionados con el Comercio (ADPIC). En el primero el TJUE precisó que la competencia interna exclusiva puede proyectarse hacia el exterior únicamente si la celebración del acuerdo internacional es necesario para realizar los objetivos del Tratado y no puedan alcanzarse mediante el establecimiento de normas autónomas. En el caso que suscitó el dictamen (el Acuerdo del GATS), el TJUE entiende que la competencia comunitaria exterior no tenía que ser exclusiva para conseguir la libre prestación de servicios y el derecho de establecimiento. Por otro lado, en el segundo, el TJUE consideró que la armonización de la propiedad intelectual no requiere la conclusión de acuerdos con terceros Estados. Por consiguiente, ambos tratados debían celebrarse como acuerdos mixtos, con competencia exterior compartida de la Unión y de sus Estados miembros.\footnote{%
    Más adelante, en el Dictamen 2/94, el TJUE consideró que su construcción sobre las competencias implícitas no confería, a las Instituciones, poderes para concluir el Convenio de Roma de 1950 (derechos humanos y de las libertades fundamentales9, al carecer las Instituciones de una competencia general en materia de protección de los derechos humanos, que sólo una modificación del Derecho originario podía conferirle.
}

En definitiva, la construcción jurisprudencial realizada por el TJUE conduce a la afirmación de la existencia de un \textbf{paralelismo entre las competencias exteriores e internas de la Unión}. En cualquiera de los casos, parece evidente que son, al menos, tres las consecuencias positivas de esa construcción del TJUE: se han ampliado considerablemente las competencias exteriores de la Unión, se ha logrado en cierta medida mitigar la rigidez del principio de atribución de competencias y, además, se ha introducido una cierta racionalidad en un sistema que, regido por dicho principio, presentaba serias deficiencias en la articulación jurídica de su dimensión exterior.

\paragraph{Competencias derivadas del recurso al art. 352 TFUE}
La última fuente del que pueden extenderse las competencias exteriores es el conocido recurso al art. 352 TFUE, que se encuentra avalado jurisprudencialmente cuando afirma que esta norma \textit{permite al Consejo adoptar todas las disposiciones apropiadas, también en el ámbito de las relaciones exteriores} (Asunto AETR).\footnote{%
    La utilización de esta norma del Tratado también conoce un valioso impulso político en la Declaración de la Cumbre de Jefes de Estado y de Gobierno de París de 1972 donde se menciona expresamente la posibilidad de acudir a esta disposición. Desde entonces, el hoy artículo 352 ha permitido una ampliación material de las competencias exteriores de la Unión actuando como fundamento jurídico exclusivo o como fundamento complementario de otra norma del Tratado, esencialmente en materia de acuerdos internacionales. El art. 352 TFUE reza lo siguiente: 
\begin{lnum}
    \item \textit{Cuando se considere necesaria una acción de la Unión en el ámbito de las políticas definidas en los Tratados para alcanzar uno de los objetivos fijados por éstos, sin que se hayan previsto en ellos los poderes de actuación necesarios a tal efecto, el Consejo adoptará las disposiciones adecuadas por unanimidad, a propuesta de la Comisión y previa aprobación del Parlamento Europeo. Cuando el Consejo adopte dichas disposiciones con arreglo a un procedimiento legislativo especial, se pronunciará también por unanimidad, a propuesta de la Comisión y previa aprobación del Parlamento Europeo}.
    \item [...]
    \item \textit{Las medidas basadas en el presente artículo no podrán conllevar armonización alguna de las disposiciones legales y reglamentarias de los Estados miembros cuando los Tratados excluyan dicha armonización}.
    \item \textit{El presente artículo no podrá servir de base para alcanzar objetivos del ámbito de la PESC y todo acto adoptado de conformidad con el presente artículo respetará los límites fijados en el párrafo segundo del artículo 40 del TUE}.
\end{lnum}
}

Se utiliza como fundamento jurídico exclusivo en materia convencional básicamente en \textbf{tres supuestos}: cuando no se puede utilizar ninguna otra competencia exterior de la Unión, cuando dicho artículo ha sido la base para la adopción de las normas internas, y cuando \textit{se utiliza junto al art. $217$ TFUE para establecer medidas provisionales y está considerado como un fundamento de carácter sustitutivo a título provisional} (RAUX, 1990: 12). 

Éste se ha empleado como fundamento jurídico complementario de muchos acuerdos, especialmente  en cuatro categorías: (i) los \textbf{acuerdos comerciales preferenciales o de cooperación comercial, económica, financiera o técnica}, que se sitúa en el amplio espacio que va desde el acuerdo meramente comercial al acuerdo de asociación;\footnote{%
    La utilización del artículo 352 junto al 207, que es el fundamento jurídico principal, se justificaba por varias razones: porque el acuerdo proyectado preveía acciones que excedían de los poderes atribuidos en el marco del artículo 207; porque se pretendía que fuese la antesala para la conclusión futura de acuerdos de asociación; porque no se querían establecer vínculos meramente comerciales ni tampoco los de carácter privilegiado que comportan los acuerdos de asociación o, incluso, por motivos meramente políticos. Este ámbito de aplicación del artículo 352 ha perdido vigor en la medida en que los ámbitos competenciales del artículo 207, tras su modificación por el Tratado de Niza, se han ampliado (comercio de servicios y aspectos comerciales de la propiedad intelectual), o se trata de los acuerdos de cooperación que encuentran desde el Tratado de Amsterdam su acomodo competencial en el artículo 211 y, naturalmente, los acuerdos sobre cooperación económica, financiera y técnica que tienen su fundamento jurídico en el artículo 212 introducido por el Tratado de Niza.
} (ii) como complemento de los arts. $186$ y $191$ en materias de investigación y desarrollo tecnológico y de medio ambiente, aunque sólo en circunstancias muy concretas; (iii) como fundamento jurídico complementario del art. $220$ relativo al establecimiento de relaciones con organizaciones internacionales y, en particular, con la ONU y sus organismos especializados, el Consejo de Europa, la OSCE y la OCDE cuando se trate de articular esas relaciones mediante acuerdos internacionales; y, (iv) como fundamento complementario de un acuerdo concluido en el ejercicio de competencias exteriores implícitas, siempre y cuando la norma en cuestión no prevea los poderes requeridos para realizar una acción necesaria para la consecución de los objetivos de la Unión (RAUX, 1990: 11-12).

En definitiva, contando con un conjunto de garantías sobre su utilización, se define como un instrumento sumamente valioso para moderar la rigidez y el estatismo propios del principio de atribución de competencias y, sobre todo, para permitir una extensión material de las competencias internas y exteriores de la Unión.\footnote{%
    El conjunto de garantías al que nos referimos es carácter juridiccional. El Consejo venía utilizando ampliamente el artículo, llegando incluso hasta el extremo de fundamentar en esa disposición determinadas actuaciones para las que contaba con una competencia y con un fundamento jurídico en virtud de otras normas del Tratado. Esta tendencia fue corregida por vía jurisprudencial, en su sentencia del Asunto 45/86, cuando afirma que \textit{sólo está justificado recurrir a este artículo como fundamento jurídico de un acto, cuando ninguna otra disposición del Tratado confiera a las Instituciones comunitarias la competencia necesaria para adoptar dicho acto}. En pronunciamientos posteriores, el TJUE confirma y precisa esta interpretación que reduce los supuestos de utilización, circunscribiéndolos a su función originaria en los términos indicados por dicha disposición.
} Además, por norma general, las competencias conferidas mediante el artículo 352 acaban por formalizarse jurídicamente en las sucesivas revisiones de los Tratados, como sucedió con el hoy artículo 212 TFUE introducido por el Tratado de Niza.

\subsection{La Política Exterior y de Seguridad Común de la Unión: Naturaleza y alcance}
Expuestas las competencias de la Unión en el ámbito de la Acción Exterior, debemos presentar la naturaleza, alcance, objetivos y competencias del \textit{proyecto comunitario} (por decirlo de alguna manera), en virtud de lo dispuesto en el Título V, Capítulo 2º del TUE, que se reduce a la conocida: Política Exterior y de Seguridad Común de la Unión Europea (PESC, en adelante), tardíamente incorporada a las relaciones exteriores comunitarias y articulada bajo la fórmula de \textbf{cooperación intergubernamental}. 

\subsubsection{Tratado de Lisboa (2007): Articulación jurídica actual}
Como decimos, siendo regulada por el TUE (en su Título V) y no por el TFUE, como expresión de su especificidad, y manteniendo unos procedimientos particulares propios y excluyentes con carácter general de la competencia del TJUE. Ambiciosa en su denominación y poco acabada aún en cuanto a su alcance y contenido, la Política Exterior y de Seguridad Común, es el fruto de una evolución iniciada en la década de los años setenta con la constitución del modelo de Cooperación Política Europea.

\paragraph{Peculiar configuración dentro de la Unión: Objetivos, principios y obligaciones}
Por lo que respecta a los objetivos, será el Consejo Europeo quien tenga atribuida la función de establecer los intereses estratégicos, los objetivos y definir las orientaciones generales de la PESC (art. 26.1 TUE), sin perder de visto los principios y objetivos generales que hemos presentado al principio.

En cuanto a las \textbf{obligaciones} (art. $24$ TUE), se pueden detraer las obligaciones generales de los Estados miembros en este ámbito: (i) «apoyo activo y sin reservas» a la PESC; (ii) «principio de lealtad y solidaridad mutua»; (iii) actuación conjunta para reforzar su solidaridad mutua; (iv) deber de abstención de acciones individuales contrarias a la Unión; y, (v) deber de abstención de acciones que puedan perjudicar o mermar la eficacia de la acción de la UE como fuerza de cohesión en las relaciones internacionales. Además, estas obligaciones generales se complementan con otras establecidas de carácter más concreto de los Estados miembros (ccoperación, coordinación, información, solidaridad, abstención de compromisos, financiación, etc.).\footnote{%
    Como es bien sabido, toda obligación, para que pueda calificarse jurídicamente como tal, debe incardinarse en un ordenamiento jurídico que garantice su eficacia y cumplimiento y la provea de mecanismos para su precisión e interpretación. Las peculiaridades de la PESC dentro del sistema jurídico de la Unión articulado en el TFUE, la exclusión de actos legislativos en esta materia, la exclusión de la competencia del TJ en las condiciones ya comentadas y el acusado carácter intergubernamental de su sistema decisorio y de aplicación sitúan a la PESC en un espacio jurídico muy peculiar, no muy alejado del sistema jurídico internacional o de organizaciones internacionales de cooperación. Esta situación no parece la más apropiada para un modelo como el diseñado en el marco de la Unión Europea, del que se presume una vocación federal o, al menos, de progresiva unificación.
}

\paragraph{Estructura jurídico-institucional: Un largo proceso inacabado}
La singularidad de la PESC, evidentmente, se proyecta en su configuración institucional. El TUE (art.13) establece la existencia de un \textbf{marco institucional único}, que constituye uno de los factores más significativos y determinantes de la estructura jurídico-funcional de la PESC. Pero, en realidad, no existe un auténtico marco institucional único. Brevemente, y habida cuenta de la diversidad de procedimientos y condiciones que las Instituciones tienen en la PESC, se consagra una diversidad funcional que lleva a considerar que el TUE se limita a instaurar \textit{una unidad orgánica y una cierta dualidad funcional} o un sistema de \textit{geometría institucional variable} y no un marco institucional unitario. Los actores más importantes, son sin duda:\footnote{
Aún más complejo resulta el sistema de representación. El artículo 15.6 TUE determina que el Presidente del Consejo Europeo asumirá la representación exterior de la Unión en los asuntos de política exterior y de seguridad común, sin perjuicio de las atribuciones del Alto Representante de la Unión para Asuntos Exteriores y Política de Seguridad. Por otro lado, se afirma que el Alto Representante representará a la Unión en las materias concernientes a la PESC, dirigirá el diálogo político con terceros en nombre de la Unión y expresará la posición de la Unión en las organizaciones internacionales y en las conferencias internacionales. Si todala modificación de Lisboa se hizo para que Europa hablara con una sola voz, por lo pronto llevamos dos. No será fácil distinguirlo. 
}
\begin{la}
    \item \textbf{Consejo Europeo}. Tiene atribuidas unas funciones excepcionalmente importantes: es el marco en el que los Estados miembros, a través de informaciones y consultas mutuas, se conciertan sobre cualquier cuestión de política exterior y de seguridad que revista un interés general con objeto de acordar un enfoque común (art. 32 TUE). Además, le corresponde adoptar decisiones relativas a los intereses y objetivos estratégicos, por él definido.\footnote{%
        Deben resaltarse tres funciones de especial relevancia atribuidas al Consejo Europeo, de carácter constitutivo de primer nivel: (i) la decisión de introducir una defensa común, decisión unánime del Consejo Europeo que deberá recomendar a los Estados miembros para que la adopten de conformidad con sus respectivas normasconstitucionales (art. 42.2 TUE); y, (ii) su capacidad para decidir por unanimidad que el Consejo pueda extender el ámbito de decisiones por mayoría más allá de lo previsto en el artículo 31.2 TUE (art. 31.3 TUE).
    }
    \item \textbf{Consejo}. Pieza clave en torno a la que pivota todo el sistema de PESC. Sus amplias y determinantes funciones pueden sistematizarse del modo siguiente: poderes decisorios e iniciativa (enfoque común, elaboración de la PESC, adopción y aplicación de decisiones, conclusión de acuerdos, medidas restrictivas, fijación del presupuesto, intermediario del PE, organización y funcionamiento del Servicio Exterior, etc.); poderes de control (coherencia, principios y consecución de éstos, control de aplicación y fiscaliza las acciones); poderes de ejecución (adopción de las decisiones sobre apliación de la PESC y PCSD); y,
    \item \textbf{Alto Representante de la Unión para Asuntos Exteriores y Política de Seguridad} (Vicepresidente de la Comisión). Pieza esencial en la competencia PESC, al Alto Representante corresponde velar por la coherencia y eficacia de la acción de la Unión, contribuir con sus propuestas a la elaboración de la PESC y a la ejecución de las decisiones, representar a la Unión en PESC y dirigir el diálogo político con terceros, organizar la coordinación de los Estados miembros en organizaciones y conferencias internacionales, y mantener las relaciones con el Parlamento Europeo. El Alto Representante cuenta para el desarrollo de estas competencias con un Servicio Europeo de Acción Exterior (SEAE, art. 27.3 TUE) que, si bien el SEAE puede parecer la materialización de un proyecto de \textit{diplomacia común europea}, su establecimiento no afecta a las bases jurídicas, responsabilidades y competencias de los Estados miembros en sus políticas exteriores.
\end{la}

Como es natural, esta complejidad se proyecta también sobre el sistema de adopción de decisiones, que a duras penas consigue abandonar la regla de la unanimidad propia de esta estructura de cooperación (modelo intergubernamental). Se pueden detectar tres tipos de reglas: (i) la unanimidad (art. 31.1 TUE); (ii) la mayoría cualificada; y, (iii) la mayoría de los miembros que componen el Consejo, aplicable a las cuestiones de procedimiento. (Es importante retener que el Consejo Europeo podrá, por unanimidad, ampliar el número de cuestiones que se rigen por la mayoría cualificada según establece el artículo 31.3 TUE).

\begin{speial}{Abstención constructivo: Una singularidad Europea (PESC)}
    Mención aparte requiere el sistema que se ha dado en llamar, no muy afortunadamente, de abstención constructiva, previsto en el apartado 1 del artículo 31 TUE. 
    
    Se trata de un mecanismo aplicableen los casos en que rige la regla general de adopción de decisiones del Consejo, es decir, la unanimidad. Es un procedimiento que pretende una cierta suavización de la regla de la unanimidad, permitiendo que las abstenciones de Estados presentes o representados no impidan la adopción de la decisión. Esto, que sería lo realmente constructivo, permitiría, pues, el nacimiento de la decisión y el despliegue de sus efectos para la totalidad de los Estados miembros, incluido, naturalmente, el o los que se hayan abstenido. Pero es ahí donde entra en juego el mecanismo previsto, según el cual uno o varios Estados (siempre que no superen más de un tercio de los Estados miembros que reúnan un tercio de la población de la Unión, en cuyo caso no habrá decisión) pueden, en el momento de la abstención, formular una declaración formal que excluye los efectos para el Estado o Estados que la formulen, admitiendo la obligatoriedad de la decisión para la Unión. Se quiere ver en ello un medio de facilitar el nacimiento de las decisiones en PESC, en la medida en que promueve la abstención que no impide esas decisiones. De ahí su caracterización como abstención constructiva y la derivación para los Estados que la ejerzan de la obligación de abstenerse de cualquier acción que pudiera obstaculizar o impedir la acción de la Unión. 
    
    Este efecto llevó a algunos a equiparar, erróneamente, este mecanismo con el de cooperación reforzada. \textbf{No deben confundirse}. Lógico: ni han sido ni son la misma cosa. Este mecanismo de la abstención constructiva incide particularmente en la obligación de los Estados implicados en la decisión de respetar la exclusión de los Estados que hayan hecho esta declaración. Si bien se observa, este mecanismo de «flexibilización de obligaciones» puede ser visto como positivo, pero también como un sistema que permite la fácil exclusión voluntaria (o inducida) de las obligaciones de PESC, afectando así negativamente a la generalidad que debiera presidir la actuación en PESC y convertirse en un instrumento más al servicio de una diversificación de obligaciones poco acorde con el principio de unidad y cohesión que se predica de la PESC.

    Todo esto tiene un sentido. El papel de los Estados miembros es fundamental en la PESC. La coordinación de su acción en organizaciones y conferencias internacionales, cuando se exija una acción operativa, las decisiones, que serán vinculantes, fijarán los objetivos, el alcance, los medios que haya que facilitar a la Unión, las condiciones de su ejecución y, en caso necesario, su duración. Las misiones diplomáticas y consulares de los Estados miembros están directamente concernidas en la aplicación de la PESC, sea con carácter general (art. 35 TUE), sea en la formulación y puesta en práctica de «enfoques comunes» (art. 32 TUE). Más particularidades ofrece la ejecución de la PCSD. La regla general es que los Estados miembros pondrán a disposición de la Unión, a efectos de la aplicación de la política común de seguridad y defensa, capacidades civiles y militares para contribuir a los objetivos definidos por el Consejo (art. 42.3 TUE).
\end{speial}

Un ámbito que, en cualquier caso, no está a la altura de la unión conseguida en los ámbitos económicos. Es lugar común apreciar las deficiencias y hasta la ineficacia del sistema de la PESC en la UE. Sin embargo, igual de común es la apreciación sobre su necesidad, lo que mantiene a la PESC entre las prioridades de la construcción europea. En cualquier caso, la articulación jurídica de la PESC, fruto de muy complicados equilibrios y sometida de manera permanente a cambios políticos internos y externos que obedecen a complejas estrategias, no es un tema fácil de comprender ni, en consecuencia, de exponer sin excesivas simplificaciones. 

En definitiva, la concepción, gestión y ejecución de la PESC se encuentran presididas por una tensión permanente entre el modelo intergubernamental y el modelo integrado, con prioridad de lo intergubernamental, construido fundamentalmente sobre el Consejo, el AR y los propios Estados miembros. Sin embargo, todo esto no ha supuesto un obstáculo para que en el marco de la PESC, partiendo de un valioso y renovado acervo acumulado, de hayan adoptado un volumen de actos (de diversa naturaleza) realmente extraordinario. Muchos de ellos revisten de una incidencia económica muy importante en tanto son parte integral de la Política Exterior de la Unión Europea que, sin naturaleza exclusivamente económica puesta que se encuentran fuera del ámbito de la Política Comercial Común y los analizados en el apartado precedente, merece la pena analizar.\footnote{%
    No entraremos a valorar la financiación de la PESC, que también tiene incidencia económica. Brevemente, en materia de \textbf{financiación} el artículo 41 TUE dispone dos reglas esenciales: (i) los gastos administrativos que la aplicación de la PESC ocasione a las instituciones «correrán a cargo del presupuesto de la Unión» siempre; y, (ii) los gastos operativos derivados de la aplicación de esta política también correrán a cargo del presupuesto de la Unión, excepto los relativos a las operaciones que tengan repercusiones en el ámbito militar o de la defensa (que tienen un régimen especial) y en los casos en que el Consejo, decida, por unanimidad, otra cosa. Cuando los gastos no corran a cargo del presupuesto de la Unión, correrán a cargo de los Estados miembros con arreglo a una clave de reparto basada en el producto nacional bruto, a menos que el Consejo decida por unanimidad otra cosa. En el caso de los gastos militares o de defensa, los Estados que hayan ejercitado la «abstención constructiva» del artículo 31.1 TUE no estarán obligados a soportar estos gastos.
}

\subsubsection{Ámbitos materiales de acción (negativo/positivo): ¿Regímenes competenciales?}
La determinación de los ámbitos materiales de la PESC se encuentra presidida por una cierta ambigüedad: no toda la política exterior es objeto de la PESC, ni ésta se identifica necesariamente con la política exterior propia de los Estados. La determinación del objeto es hoy fruto de una doble operación.

\paragraph{Sentido negativo: Política, no Acción Exterior}
En ningún caso forman parte del ámbito material de la PESC los ámbitos de Acción Exterior que han sido atribuidos a la Unión Europea como políticas de la acción exterior de las competencias contempladas en los artículos 3 a 6 TFUE (las relaciones exteriores que hemos estudiado en los dos capítulos anteriores). La PESC se configura como una competencia diferente regulada sólo en el TUE. Dicho de otra manera, los títulos competenciales PESC (TUE) y los situados en el TFUE relativos a la acción exterior que no son PESC (herederos de los de las CCEE), y al que se ha asimilado el del ELSJ, son sustancialmente distintos: la \textbf{acción exterior comunitaria} (TFUE) proviene de la acción exterior esencialmente referida a \textbf{ámbitos económicos a los que se ha sumado el ELSJ} y que \textbf{la acción exterior PESC es la esencialmente política};

\paragraph{Sentido positivo: }
La segunda operación nos conduce a la identificación positiva de los contenidos materiales dela PESC: sólo cuando se produce un consenso de los Estados sobre la existencia de asuntos o cuestiones de política exterior que ofrecen un interés general o común, ese asunto o cuestión pasa a formar parte de la PESC.\footnote{%
    Aun aasí, es evidente que en el TUE se han identificado algunos ámbitos generales (art. 21.2 TUE) y algunos más específicos (así, las acciones en gestión de crisis previstas en el art. 42, cuyo alcance se concreta en el art. 43 TUE), pero lo que establece son esencialmente procedimientos para la identificación concreta de los ámbitos materiales, a través de ese consenso, y atribuye al Consejo Europeo esa función cuando dispone que: \textit{El Consejo Europeo determinará los intereses estratégicos de la Unión, fijará los objetivos y definirá las orientaciones generales de la política exterior y de seguridad común, incluidos los asuntos que tengan repercusiones en el ámbito de la defensa, adoptando para ello las decisiones que resulten necesarias} (art. 26.1 TUE).
}

Por tanto, la cuestión del ámbito material de la PESC podría resolverse afirmando que cubre el conjunto de la política exterior plenamente en sus aspectos políticos y económicos, y de manera más condicionada aquellos aspectos de la seguridad que tengan relación con los aspectos militares y de la defensa.\footnote{%
    Con relación a estos últimos, divididos en dos vertientes, la \textit{política de defensa común} y la \textit{defensa común} (ésta no incorporada). Así, en el primer caso, la política de defensa común está incorporada por el propio TUE (art. 42) aunque exige un desarrollo y una definición especial. El segundo, se limita a prever su posible incorporación, para lo que exige una decisión unánime del Consejo Europeo recomendándolo y la aceptación de los Estados miembros de conformidad con sus respectivas normas constitucionales (art. 42.2 TUE).
}

Los objetivos materiales de la PESC ofrecen un indicador pero no una solución a la concreción de su ámbito material. La tradicional pretensión de que el TUE concretase y predefiniese los ámbitos materiales no ha conseguido acuerdo. Sin embargo, podríamos aproximarlos desde un ámbito descriptivo.\footnote{%
    De ahí la atribución al Consejo Europeo de la función de fijar los objetivos estratégicos y definir las orientaciones generales de la PESC sobre cuya base el Consejo toma las decisiones necesarias para definir y ejecutar la política exterior y de seguridad común (art. 26 TUE). El juego de los principios de subsidiariedad y proporcionalidad no parece comportar una seria restricción en este campo. Lo que conduce a pensar que, existiendo una confluencia de voluntades de los Estados miembros, cualquier cuestión puede ser objeto material de la PESC.
} La PESC parte ahora de un valioso \textit{acervo} acumulado en estos años anteriores. La acción en materia de PESC se distribuye por temas (p. ej., la lucha contra el terrorismo), por proximidad geográfica (Política Europea de Vecindad), por zonas geográficas (determinadas lógicamente según zonas de interés en PESC y, dentro de ellas, por países) o por zonas de conflicto, por su incardinación en problemas de seguridad y defensa (gestión de crisis, prevención de conflictos, desarme y no proliferación, etc.). Junto a ello, a veces se aborda la participación PESC en función de grandes cuestiones temáticas (así, Corte Penal Internacional) o foros mundiales (como las Naciones Unidas), entre otros.


La concepción general de la PESC responde esencialmente al rechazo de las propuestas de «comunitarización» de la política exterior europea (art. 24 TUE): \textit{La política exterior y de seguridad común se regirá por reglas y procedimientos específicos. La definirán y aplicarán el Consejo Europeo y el Consejo, que deberán pronunciarse por unanimidad salvo cuando los Tratados dispongan otra cosa. Queda excluida la adopción de actos legislativos}.\footnote{%
    El Tribunal de Justicia de la Unión Europea no tendrá competencia respecto de estas disposiciones, con la salvedad de su competencia para controlar el respeto del artículo 40 del presente Tratado y para controlar la legalidad de determinadas decisiones contempladas en el párrafo segundo del artículo 275 del Tratado de Funcionamiento de la Unión Europea, aunque no afecta a la exclusión general. No obstante, se reconoce que (art. 37 TUE) \textit{La Unión podrá celebrar acuerdos con uno o varios Estados u organizaciones internacionales en los ámbitos comprendidos en el presente capítulo}, que concluidos por vía del artículo 218.3 TFUE, se generan actos con evidentes efectos jurídicos que el TJUE ha entendido sometidos a su competencia. Paradigmático es el caso en que el TJUE anuló la Decisión 2011/640/PESC del Consejo, relativa a la firma y la celebración del Acuerdo entre la Unión Europea y la República de Mauricio sobre las condiciones de entrega, por la fuerza naval dirigida por la Unión Europea a la República de Mauricio, de sospechosos de piratería y de los bienes incautados relacionados, y sobre las condiciones de trato de tales sospechosos después de su entrega.

Además, se incluye un artículo que confirma la dualidad de concepciones para la PESC y otras competencias de la Unión, asegurando la \textbf{no contaminación} de uno y otro espacio.
} 

Como vemos, se trata de una estructura, claramente diferenciada del sistema general, que se concibe en el marco de la Unión Europea como una competencia que abarcará todos los ámbitos de la política exterior y todas las cuestiones relativas a la seguridad de la Unión, incluida la definición progresiva de una política común de defensa que podrá conducir a una defensa común.\footnote{%
    Esta diferenciación hace que, pese a la insistencia en presentar la acción exterior como algo unitario, persista el gran problema de las posibles contradicciones entre unas actuaciones de la Unión en el ámbito general de su acción exterior y la específica de la PESC. No es de extrañar que asegurar la coherencia de la actuación de todos los instrumentos de acción exterior sea objeto de una preocupación fundamental. En este contexto, el Consejo y la Comisión tienen atribuida la responsabilidad de garantizar, en sus respectivas competencias, dicha coherencia, a cuyo fin les impone el TUE una obligación de «cooperación»
} En este sentido, la Unión dirigirá la PESC (art. 25 TUE): mediante la definición de sus orientaciones generales, con la adopción de decisiones de una triple naturaleza (acciones, posiciones y modalidades de ejecución de éstas) y fortaleciendo la cooperación sistemática entre los Estados miembros para llevar a cabo sus políticas.





\clearpage
\section{Acción Exterior de la Unión: Relaciones económicas exteriores}
\puntosclave{
    
}


\subsection{Acción Exterior}



\subsubsection{Políticas sobre control de fronteras, asilo e inmigración: Incidencia económica}
Las Políticas sobre control de fronteras, asilo e inmigración se conocen también como \textbf{Política Común de movimiento de trabajadores}. Como hemos visto, esto ha sido fruto de una evolución histórica muy largo y, hoy en día, constituye uno de los elementos más característicos de la Unión.

\paragraph{Resultados: Trabajo extranjero en la UE}
Como sabemos, en la migración influye una combinación de factores económicos, medioambientales, políticos y sociales: tanto en el país de origen del inmigrante (factores de empuje) como en el país de destino (factores de atracción). La relativa prosperidad económica y la estabilidad política de la UE han ejercido un considerable efecto de atracción sobre los inmigrantes. En los países de destino, la migración internacional puede utilizarse como herramienta para hacer frente a carencias especificas del mercado laboral. Sin embargo, es casi seguro que la inmigración por sí sola no invertirá la tendencia actual de envejecimiento de la población que se experimenta en muchas partes de la UE.

Tras un descenso significativo del número de migrantes en 2020 debido a la pandemia de COVID- 19 la migración internacional se recuperó en 2021 (inmigración neta positiva de 960.000 personas) y los datos disponibles sugieren que siguió creciendo en 2022. En 202 1 los ciudadanos de países no U E representaban el 5,3\% de la población total de la U E (los ciudadanos residentes nacidos fuera de la UE alcanzan el 8,4\% de la población total). Ese año entraron 2,3 millones (aumento del 18\% respecto a 2020). Los motivos principales de resistencia fueron por trabajo y razones familiares (20\% y 36\% de los permisos de residencia expedidos). Destaca el importante aumento relativo de los permisos por trabajo. Los inmigrantes representaban un 4,7\% de la población activa total. La tasa de empleo de ciudadanos de países no pertenecientes a la UE fue del 59, 1\% (frente al 74\% de ciudadanos de la UE). Aparecen sobrerrepresentados en los sectores de hostelería y restauración, actividades  administrativas y servicios auxiliares, trabajo doméstico y construcción, asociados generalmente a una menor cualificación. 

En la UE, Polonia expidió el mayor número de primeros permisos de residencia a ciudadanos de fuera de la UE (33\% del total) seguida de España (13\%), Francia (10\%), Italia (9\%), Alemania (6\%) y Países Bajos (4\%).

Si se tiene en cuenta a los inmigrantes procedentes de otros países UE, Alemania registró el mayor número total de inmigrantes (874.400) en 2021, seguida de España (528.900), Francia (336.400) e Italia (3 18.400). Por su parte, Croacia, Grecia, Letonia y Rumanía presentaron un mayor número de emigrantes en 202 1 . Por otro lado, Luxemburgo presentó en 2021 la tasa más alta de inmigración por cada 1.000 habitantes.

La información sobre ciudadanía se ha utilizado a menudo para estudiar a los inmigrantes de origen extranjero. Sin embargo, dado que la nacionalidad puede cambiar a lo largo de la vida de una persona, también es útil analizar la información por país de nacimiento. Los países de origen de los residentes nacidos en el extranjero eran en 202 1 por orden de importancia: Singapur, Australia, Nueva Zelanda, Suiza, Canadá, Islandia, Noruega, Estados Unidos, Reino Unido, Ucrania. Chile y Rusia.

En 2021 inmigraron a la UE ligeramente más hombres que mujeres (55% frente a 45%). El Estado miembro que comunicó la mayor proporción de inmigrantes varones fue Croacia (73%); por el contrario, la mayor proporción de inmigrantes mujeres se registró en Chipre (54 %).

A I de enero de 2022, la edad media de la población nacional de la UE era de 45,4 años, mientras que la de los no nacionales residentes en la UE era de 36,6 años. Es probable que la mayoría de los inmigrantes que residen en la UE se queden a medio plazo, por lo que, además de un empleo, el apoyo social, como el acceso a la vivienda, la asistencia sanitaria, la educación y la ayuda a los niños, es igualmente importante para su integración.

Entre los principales retos para la integración de estas personas se encuentra: el acceso al mercado laboral; reducir la brecha de acceso al empleo de las mujeres, mejorar la educación y la inclusión social. En todo caso, abordar los retos de la integración de los inmigrantes es un proceso multidimensional y requiere coordinar acciones políticas en distintos ámbitos. Desde la presentación del Plan de· Acción sobre la Integración de Nacionales de Terceros Países, en junio de 2016, la Comisión ha trabajado para reforzar el enfoque común en todos los ámbitos políticos e implicar a todos los agentes pertinentes.\footnote{%
    La Comisión apoya la integración temprana en el mercado laboral, los servicios públicos de empleo, el acceso a la educación y la formación. Puso en marcha en 20 1 7 la herramienta de perfil de competencias de la UE para nacionales de terceros países. una herramienta multilingüe en línea para ayudar a identificar y cartografiar las competencias y cualificaciones. Asimismo, supervisa la evolución de las políticas de empleo e inclusión social de estos grupos vulnerables a través del Semestre Europeo que aplica la Estrategia Europa 2020 y financia una serie de medidas para la integración de los inmigrantes a través del Fondo Social Europeo. el Fondo de Ayuda Europea para los Más Necesitados y el Programa de Empleo eInnovación Social. Por último, fomenta y apoya las políticas basadas en pruebas. el aprendizaje mutuo, l
}

Por último, cabe mencionar la confrontación existente entre los PEO y las economías avanzadas en relación al modo 4 del comercio de servicios (ver Anexo 7).

En virtud del Acuerdo General sobre el Comercio de Servicios (AGCS o GATS en inglés), los servicios pueden suministrarse a nivel internacional de cuatro formas distintas, conocidas como "modos de suministro". El modo 4 se refiere a los servicios suministrados por personas de un Miembro de la OMC mediante su presencia en el territorio de otro M iembro (presencia de personas físicas). Abarca a los empleados de empresas de servicios y a los proveedores de servicios por cuenta propia, excluyendo las medidas en materia de ciudadanía, residencia o empleo con carácter permanente y las medidas que afecten a personas que traten de acceder al mercado de trabajo de otro Miembro.

Esta exclusión que protege a los países desarrollados de las demandas de liberalización de la migración es objeto de controversia y dura crítica por parte de los PEDs, tradicionalmente emisores de mano de obra poco cualificada. Uno de los problemas subyacentes es la confusión existente entre movilidad laboral y movilidad del servicio. Además, el AGCS no define qué se entiende por "acceso al mercado de trabajo" que en última instancia depende de las normativas nacionales. M ientras que, en principio, el AGCS permite la circulación de todas las categorías de personas físicas, en la práctica los Estados miembros han liberalizado la circulación de categorías muy específicas de personas físicas. Muy pocos miembros de la OMC han asumido compromisos significativos en el marco del Modo 4 y los compromisos existentes se refieren principalmente a la circulación de personas trasladadas dentro de una misma empresa (en la mayoría de los casos, ejecutivos, directivos y especialistas), visitantes de negocios y trabajadores autónomos altamente cualificados. Ese bajo incentivo se debe también a la obligación de aplicar el principio de Nación Mas Favorecida. Además, los compromisos del Modo 4 no suelen ser vinculantes para el statu quo, sino que a menudo reflejan condiciones de entrada más estrictas que las que se conceden en la práctica. Teniendo en cuenta que los compromisos del AGCS se contrajeron a mediados de los años noventa, es probable que- la diferencia entre las condiciones de entrada vinculantes y las aplicadas se haya acentuado.

De hecho, el comercio bajo modo 4 es muy residual: en 2019 el modo 4 representó el 2,9\% del comercio mundial de servicios, lo que también puede estar reflejando un acceso bastante limitado a los mercados.

Los compromisos otorgan derechos de entrada y estancia únicamente a determinadas categorías de personas, en particular a las personas ligadas al establecimiento de una presencia comercial (por ejemplo, el personal transferido dentro de la misma empresa) y las personas altamente cualificadas (como directivos, ejecutivos o especialistas). Además, es frecuente encontrar otras restricciones, como contingentes, "pruebas de necesidades económicas" o "pruebas del mercado de trabajo" (pruebas que supeditan el acceso al cumplimiento de determinados criterios económicos), condiciones relativas al empleo anterior y requisitos en materia de residencia y capacitación.

Cada vez es más necesario transferir eficazmente conocimientos especializados y mano de obra internacionalmente, ya sea mediante el traslado temporal de especialistas o con fines de servicios contractuales. Aunque estas transacciones no sustituyen a la inmigración tradicional por motivos laborales, pueden aumentar los beneficios y aportar ingresos adicionales tanto a los Estados de origen como a los de acogida. Una consideración más práctica es que sería poco realista esperar que los países de la OCDE abran sus mercados' a la movilidad laboral mediante un acuerdo multilateral jurídicamente vinculante. Quizá sería más productivo pedir compromisos más profundos que sigan siendo vinculantes el ámbito del comercio de servicios. La liberalización de este modo de suministro puede jugar también un papel importante como baza negociadora para conseguir otros objetivos.

\paragraph{MFP 2021-2027: Migración y gestión de las fronteras (FRONTEX y FSAMI)}
Ello ha llevado a la creación de agencias (como la Agencia europea para la gestión de la cooperación operativa en las fronteras exteriores —Frontex—) y al establecimiento de mecanismos de asistencia (es el caso de los equipos de intervención rápida en fronteras) y de compensación entre EEMM, así como de instrumentos financieros europeos, donde destacan el \href{https://fondoseuropeosparaseguridad.interior.gob.es/es/fondos/fondo-de-seguridad-interior/}{Fondo de Seguridad Interior} y el \href{https://www.inclusion.gob.es/web/migraciones/fami}{Fondo de Asilo, Migración e Integración}.\footnote{%
    A pesar de los instrumentos europeos adoptados, la puesta en práctica de esta solidaridad y reparto equitativo de las cargas se ha enfrentado a muy serias resistencias de los EEMM y, a decir verdad, constituye aún un objetivo a lograr. Acontecimientos excepcionales, como la denominada crisis de los refugiados, han hecho emerger esta ausencia de voluntad política común que, unida a las deficiencias de su configuración jurídica, ha colocado las políticas de este capítulo en serio peligro de colapso.
}


\subsubsection{Relaciones con otros sujetos \textit{ex} art. 217: Asociación}
Pasamos a analizar las relaciones económicas de la UE con sus principales socios en términos geoestratégicos.

\textbf{Bloque occidental}. Son países con los que la UE comparte los valores de la democracia, los derechos humanos, el Estado de Derecho y la libertad económica y política, y tienen intereses comunes en materia de política exterior y seguridad. La estrecha colaboración y las relaciones estratégicas con éstos son una prioridad para la Unión.

\textbf{BRICS}. BRICS es el acrónimo de una asociación informal de contenido económica-comercial de las cinco economías nacionales emergentes que en la década de los 2000 eran las más prot;1etedoras del mundo y considerados en origen el paradigma de la cooperación Sur-Sur. Desde su cumbre germinal en 200 9, se han celebrado reuniones de los jefes de Estado de los BRICS en catorce ocasiones. Estas economías tienen en común una gran población (China e India por encima de los mil cien millones, Brasil y Rusia por encima de los ciento cuarenta millones), un amplio territorio (casi 38,5 millones km2) dotado de una gigantesca cantidad de recursos naturales y, unos volúmenes de PIB entre los mayores del mundo, acompañados de altas tasas de crecimiento y de incremento en la participación en el comercio mundial en los últimos años.

\paragraph{Relaciones transatlánticas: Estados Unidos y Canada}
La UE y los EE.UU. cooperan estrechamente en una serie de ámbitos de política exterior y contextos geográficos, tales como la lucha contra el terrorismo, la cooperación en materia de seguridad o la energía.

En diciembre de 2020 se publicó una nueva agenda UE- EE.UU. para el cambio mundial. en la cual se identificaron oportunidades de cooperación en relación con la pandemia mundial. el clima. el comercio, la seguridad y la democracia. Desde entonces, se han puesto en marcha varios grupos bilaterales EE.UU.-UE sobre comercio y tecnología (TTC), sobre seguridad y defensa, sobre China y el Indo-Pacífico y sobre Rusia, mientras que se ha reactivado el Consejo de Energía y, tras la agresión rusa contra Ucrania, se hacreado un grupo de trabajo específico para reducir la dependencia europea de los combustibles fósiles rusos.

De hecho, hay un gran número de cuestiones en las que se puede trabajar conjuntamente para lograr resultados mutuamente beneficiosos (por ejemplo, la reforma de la OMC) y en las que ya se han realizado importantes avances, como en el ámbito de la armonización del impuesto de sociedades, o la superación de las diferencias en cuanto a subsidios a aeronaves, aranceles sobre acero y aluminio y flujos de datos

La Unión y los EE.UU. son los principales inversores mutuos, con una inversión total estadounidense en la Unión tres veces superior a la de Asia en su totalidad. La inversión de la Unión en EE.UU. es aproximadamente ocho veces superior a la inversión de la Unión en India y China juntas. Podría afirmarse que la inversión directa bilateral, que por definición es un compromiso a largo plazo, es la fuerza motriz de las relaciones comerciales transatlánticas. Lo confirma el hecho de que el comercio entre sociedades matrices y filiales en la Unión y en los EE.UU. constituye más de una tercera parte de todo el comercio transatlántico. Según algunas estimaciones, las empresas de la Unión y de los EE.UU. que operan en el territorio del otro proporcionan empleo a más de 14 millones de personas.

Ante la entrada en vigor del Intlation Reduction Act (IRA) en agosto de 2022, se ha creado una Task Force entre la U E y EE.UU. en la que se están negociando medidas que puedan amortiguar el impacto del IRA en las empresas europeas y preservar la cooperación entre ambas economías. En cuanto a Canadá, las relaciones bilaterales con este país comenzaron en la década de 1950 por razones económicas y desde entonces han ido evolucionando hasta convertirse en una estrecha asociación estratégica. El Acuerdo de Asociación Estratégica U E-Canadá (AAE) es un acuerdo político exhaustivo que busca reforzar la cooperación bilateral en una serie de ámbitos sectoriales y de política exterior, por ejemplo, la paz y la seguridad internacionales, la lucha contra el terrorismo, la gestión de crisis, la seguridad marítima. la gobernanza mundial. la energía, el transporte. la investigación y el desaiTollo, la salud, el medio ambiente, el cambio climático y el Ártico. Partes sustanciales del Acuerdo están en vigor, con carácter provisional desde 2017.

Por otro lado, el Acuerdo Económico y Comercial Global (CETA) es el resultado de la evolución positiva de las relaciones comerciales UE-Canadá durante la última década.\footnote{%
    El CETA tamhién es el primer acuerdo económico bilateral de la Unión que incorpora un Sistema de Tribunales de Inversiones para la solución de diferencia en materia de inversión entre inversores y Estados. Daqo su carácter innovador y teniendo en cuenta que d dehate público al respecto todavía no se ha ccrrndo en muchos países, el sistema no entra en el ámbito de la aplicación provisional del CETA. Asimismo, el CETA contiene una declaración int:quívoca sobre el derecho de los gobiernos a regular con tines públicos en materia de salud pública seguridad. medio ambiente. moral pública y protecdón social y del consumidor.
} Ha mejorado considerablemente las relaciones económicas, comerciales y de inversión entre la Unión y Canadá al abrir sus mercados a los bienes, servicios e inversiones del otro, incluida la contratación pública.

\paragraph{Relaciones con Rusia}
Las relaciones entre la UE y Rusia se basan en el acuerdo de colaboración y cooperación de 1994, con un periodo de vigencia inicial de I O años, que finalmente se ha ido renovando anualmente de forma automática. En 200 3 se reforzó la cooperación entre ambos agentes mediante la creación de cuatro espacios comunes.\footnote{%
    Espacio económico, un espacio de libertad, seguridad y justicia un espacio de seguridad exterior y un espacio de investigación, educación y cultura.
} En 2008 se iniciaron negociaciones con el objetivo de alcanzar un nuevo acuerdo que contara con compromisos jurídicamente vinculantes en diversas áreas ( dialogo político, justicia, cooperación económica, energía). Adicionalmente, la UE fue una firme defensora de la Adhesión de Rusia a la OMC, que finalmente tuvo lugar en 20 12. 

Sin embargo, la relación entre la UE.y Rusia ha basculado en la última década, debido a las fuertes tensiones en la frontera entrt: ambas potencias. Primero, la incorporación de Crimea a Rusia en 20 14 llevó a una revisión de la relación bilateral de la UE con Rusia y llevó a la suspensión de. las cumbres periódicas y las negociaciones sobre el nuevo acuerdo bilateral, la exclusión de Rusia del G8 y la interrupción del proceso de adhesión de Rusia a la OCDE y a la Agencia Internacional de la Energía, ade . más de la imposición de sanciones económicas.

Segundo, la agresión militar rusa a Ucrania iniciada en 2022 y la anexión por parte de Rusia de cuatro regiones ucranianas27 han sido condenadas por la UE y sus miembros. Esto ha llevado a adoptar nuevamente un conjunto de sanciones, de una severidad sin precedentes con el fin de debilitar la economía rusa, privándola de acceso a tecnologías, recursos y mercados clave. Asimismo, la UE tanto en el marco de la OTAN, como a nivel comunitario, como a nivel de cada Estado miembro, ha optado por apoyar militar y económicamente al gobierno ucraniano.

En 2022 Rusia era el mayor proveedor de gas y petróleo de los EEMM, teniendo en muchos de ellos una posición cercana al monopolio. Dada la elevada dependencia de los combustibles fósiles rusos por parte de la UE, el empeoramiento de las relaciones entre los dos ejes ha dado lugar a que la soberanía energética se haya convertido en una cuestión de máxima prioridad. Para superar esta vulnerabilidad, la UE lanzó en mayo de 2022 REPowerEU, un plan para poner fin a las importaciones de energía fósil de Rusia mucho antes de 2030. Para ello es necesario diversificar el suministro energético de la UE, aumentar el ahorro y la eficiencia energética y acelerar la transición hacia una energía verde.

\paragraph{Relaciones con Asia (meriodional): India}
La asociación estratégica entre la Unión y la India fomenta el comercio y la cooperación económica. Cada socio cuenta con regiones con grandes diferencias en cuanto a fortaleza económica, idioma y cultura, un gran mercado y una posición geocstratégica que plantea preocupaciones en materia de política de seguridad.

La Unión es el tercero de los mayores socios comerciales de la India, con un comercio de mercancías por valor de 62.800 millones de euros en 2020 (un 11,1\% del total del comercio de la lndia) y una balanza comercial de alrededor de 900 millones de euros a favor de la India. La Unión es el mayor inversor extranjero en la India. y los flujos de inversión extranjera en el país han aumentado del 8\% al 18\% en la última década. La inversión extranjera directa en la India ascendió a un total de 75.800 millones de euros en 2019. La India se beneficia actualmente de un trato arancelario preferencial unilateral en el marco del Sistema de Preferencias Generalizadas de la Unión (SPG), que vincula las preferencias comerciales unilaterales al respeto de los derechos humanos y laborales.

En 2021 se llegó a un compromiso para reanudar las conversaciones sobre el estancamiento del acuerdo de libre comercio y de iniciar las negociaciones sobre un acuerdo de protección de las . inversiones y un acuerdo sobre indicaciones geográficas. La India espera que se pueda firmar antes de sus próximas elecciones en 2024.

La Asociación sobre Conectividad UE-lndia también se puso en marcha en 202 1 con vistas a apoyar una conectividad resiliente y sostenible tanto en la India como en terceros países. Estos avances son importantes en el contexto de las tensiones con China.

La UE y la India publicaron su Estrategia para la región indopacífica en 2021. haciendo hincapié en un orden internacional multilateral basado en no1mas de la región, en particular como demostración de fuerza de cara a China. La India se ha comprometido a garanti:zar un ciberespacio seguro y a mejorar la seguridad marítima.

\subsubsection{Acuerdo Euromediterráneo de Asociación con Israel}
El Acuerdo Euromediterráneo de Asociación con Israel, firmado en 1995 en el marco del Proceso de Barcelona y en vigor desde el año 2000, liberaliza el comercio de productos industriales y contiene disposiciones limitadas sobre productos agrícolas. Como en todos los Acuerdos de Asociación de la Unión, su artículo 2 configura el respeto de los derechos humanos y de los principios democráticos como «elemento esencial» del Acuerdo (la misma técnica jurídica, análoga a una cláusula democrática, que hemos visto incorporarse por primera vez a un ámbito distinto de los derechos humanos en el capítulo climático del acuerdo con Mercosur).

\paragraph{Revisión del cumplimiento (\textit{ex} art. 2: Conflicto en Gaza} 
A raíz del conflicto en la Franja de Gaza, el Alto Representante de la Unión presentó al Consejo de Asuntos Exteriores, el 23 de junio de 2025, una revisión que concluyó la existencia de indicios de incumplimiento por parte de Israel del artículo 2 del Acuerdo.\footnote{%
    Consejo de Asuntos Exteriores de la UE, 23 de junio de 2025; Centro de Documentación Europea, Universidad de Granada.
} En septiembre de 2025, la presidenta de la Comisión anunció una propuesta de suspensión parcial de las disposiciones comerciales del Acuerdo, junto con sanciones específicas dirigidas a determinados ministros y a colonos vinculados a episodios de violencia en Cisjordania. A la fecha de las últimas informaciones disponibles (mayo de 2026), esa propuesta seguía sin obtener la mayoría necesaria en el Consejo, bloqueada singularmente por Alemania e Italia, por lo que el Acuerdo permanece formalmente en vigor en su integridad.\footnote{%
    Amnistía Internacional, "Líneas rojas, no alfombras rojas", mayo de 2026; se trata de un asunto políticamente controvertido y en evolución activa, sobre el que distintos actores institucionales y organizaciones mantienen posiciones encontradas.
} Este caso constituye, junto con el más antiguo precedente de la jurisprudencia \textit{Brita} (2010) --sobre el origen preferencial de mercancías procedentes de los asentamientos-- y la sentencia \textit{Psagot} (2019) --sobre el etiquetado de origen de productos de los asentamientos--, uno de los ejemplos más citados en la doctrina para ilustrar la tensión entre la cláusula de derechos humanos de los Acuerdos de Asociación y su aplicación práctica efectiva.



\subsection{Política Exterior y de Seguridad Común}
La acción en materia de PESC se distribuye por temas (p. ej., la lucha contra el terrorismo), por proximidad geográfica (Política Europea de Vecindad), por zonas geográficas (determinadas lógicamente según zonas de interés en PESC y, dentro de ellas, por países) o por zonas de conflicto, por su incardinación en problemas de seguridad y defensa (gestión de crisis, prevención de conflictos, desarme y no proliferación, etc.). Junto a ello, a veces se aborda la participación PESC en función de grandes cuestiones temáticas (así, Corte Penal Internacional) o foros mundiales (como las Naciones Unidas), entre otros.

Tradicionalmente, pues, la PESC ha partido de una separación de la seguridad en tres vertientes: política, económica y militar (la defensa en sentido estricto). Sobre esa triple dimensión el TUE ha mantenido siempre que la seguridad forma parte de la PESC en sus aspectos políticos y económicos. En consecuencia, la defensa (aspectos militares de la seguridad) no. La presión política de algunos Estados miembros abrió una brecha distinguiendo entre «política de defensa común» y «defensa común». La primera ha permitido operar sobre algunos aspectos político-estratégicos de la defensa mientras que la segunda apuntaría en el sentido de dotar a la UE de un sistema dedefensa propio, asunto más complejo. La primera, la política de defensa, permitió una cierta apertura y ha tenido importantes desarrollos. La segunda, por el contrario, ha estado siempre y aún permanece cerrada hasta la obtención de una decisión unánime del Consejo Europeo y su ratificación por todos los Estados miembros según sus respectivas normas constitucionales (art. 42.2 TUE). 


\subsubsection{Proximidad geográfica: Política Europea de Vecindad}
Con la ampliación de 2004, las fronteras de la UE se modifican notablemente, lo que hace que países que antes estaban más apartados de la UE pasen a ser ahora países vecinos. En este contexto, la UE lanza la denominada "Política Europea de Vecindad" (PEV) para ofrecer a sus vecinos una relación privilegiada, basada en el compromiso mutuo con unos valores comunes (la democracia y los derechos humanos, el Estado de Derecho, la buena gobernanza, los principios de la economía de mercado y el desarrollo sostenible).\footnote{%
    La Política Europea de Vecindad. no es el primer imento de acercamiento de la UE respecto a sus vecinos mediterráneos. La colaboración existía desde 1995 a través del Pacto de Barcelona (""Estrategia de Asociación Euromt:dite1Tánea o Euromed") y. en 2008, pasó a denominarse la --unión por el Mediterráneo" y a tener instituciones propias.
} A su vez, con esta política, la UE busca fortalecer la prosperidad, la estabilidad y la seguridad en sus fronteras. La base jurídica de esta política se encuentra en el Artículo 8 del TUE, el título V del TUE (acción exterior) y los artículos 206 y 207 (comercio), y 216 a 219 (acuerdos internacionales) del Tratado de Funcionamiento de la UE (TFUE). La PEV es una política con una doble dimensión: bilateral, entre la UE y cada país socio, y multilateral, dado que dichas relaciones bilaterales son completadas con iniciativas regionales de colaboración: la Asociación Oriental y la Unión por el Mediterráneo.

En 2015 se llevó a cabo una revisión de la política y desde entonces ha conseguido movilizar un volumen importante de apoyo en cuatro esferas prioritarias de reforma:
• la buena gobernanza, la democracia, el Estado de Derecho y los derechos hum��nos, con amplios programas de reforma de la administración pública y lucha contra la corrupción, consolidación del poder judicial. ..
• el desarrollo económico sostenible, fundamental en la contribución de la VE a la estabilización de la vecindad y es crucial para fomentar la resiliencia económica de los países socios
• la seguridad;
• la migración y la movilidad, puesto que la crisis de los refugiados y la migración ilegal siguen siendo prioritarios en el programa político y son un aspecto clave de la colaboración de la UE con sus países vecinos.

Un informe sobre dicha revisión fue publicado en 2017, en el cual se destacaba la mayor flexibilidad y sensibilidad con la que estaba actuando la Unión y el uso más eficiente de los recursos.

En 2019, el Parlamento aprobó una Resolución titulada « Después de la Primavera Árabe: el camino a seguir en la región MENA», en la que, si bien reconocía que se habían producido algunos avances positivos, pedía nuevas refonnas económicas, democráticas y sociales. En 2021 se publicaron dos documentos: (i) la Comunicación «Asociación renovada con los países vecinos meridionales - Una nueva Agenda para el Medfterráneo» y, (ii) el «Plan de acompañamiento de inversiones para los vecinos meridionales», que buscaban renovar la PEV para relanzar y reforzar la asociación estratégica entre la Unión y sus vecinos meridionales.

Los países que abarca esta política pueden dividirse en dos grupos principales:

a) Europa Oriental: .Ucrania, Moldavia, Georgia, Bielorrusia, Annenia, Azerbaiyán .. En este ámbito de países existe la Asociación Oriental, cuyo objetivo es mejorar las relaciones de la Unión con la mayor parte de sus vecinos del este. Resulta importante resaltar que Ucrania, Moldavia y Georgia han pasado a· ser candidatos a la adhesión o potenciales candidatos en 2022
b) Mediterráneo Oriental y Meridional (Marruecos, Argelia, Túnez, Libia, Egipto, Israel, la Autoridad Palestina, Jordania, Siria y Líbano). La UE tiene firmados acuerdos de asociación con todos estos países, con la excepción de Libia. 

La política de vecindad se fundamenta en dos principios: la copropiedad_ y la diferenciación. La PEV no viene impuesta por parte de la UE, sino que es definida en colaboración con los países implicados. Adicionalmente, esta política no es homogénea para todos ellos, sino que se tiene en cuenta la situación específica de cada país. Es decir, pese a tratarse de una política común para un conjunto de países, la UE tiene flexibilidad para adaptar las acciones concretas a las características específicas de cada socio, algo indispensable dada la heterogeneidad de los países en cuestión. l11strume11tos

El principal instrumento de la PEV son los planes de acción bilaterales que contienen una agenda de reforrnas políticas y económicas a llevar a cabo en estos países y donde se identifican unas prioridades a corto y medio plazo (de 3 a 5 años). 

Desde el punto de vista de la financiación, la cooperación de la UE durante el periodo 2021 -2027 se enmarcará en el nuevo Instrumento de Vecindad, Cooperación al Desarrollo y Cooperación Internacional "Europa Global" (NDICI, por sus siglas en inglés):

• Se regula en el Reglamento 2021/947 que entró en vigor 9 de junio de 2021 pero se aplicó retroactivamente desde el I de enero de 2021. Supone un cambio estructural importante porque aglutina I O instrumentos de financiación previos (8 Reglamentos, 1 Decisión y 1 Fondo extrapresupuestario) en un único instrumento simplificando la arquitectura de financiación exterior de la UE. Está dotado con 79.500 millones de euros en precios de 202 1 ; es decir, un importe sustancialmente mayor al de los IPA. Además, como éstos, la dotación se ha incrementado notablemente respecto al anterior MFP (+ 1 2%). Está prevista una evaluación de este Instrumento para diciembre 2024.

• Los ocho reglamentos son los relativos af Instrumento de Cooperación al Desarrollo (ICD), el Instrumento Europeo de Vecindad (IEV), el Instrumento para la paz y estabilidad, el Instrumento europeo para la democracia y los derechos humanos, el Instrumento de partenariado de cooperación con terceros países, el Reglamento de reglas comunes, tres mecanismos de garantía que incluyen dos reglamentos (el Fondo Europeo de Desarrollo sostenible FEDS y el Fondo de Garantía de Acción Exterior) y una Decisión que regula la garantía al BEi para cubrir pérdidas de operaciones financieras del BEi en el exterior (mandato de préstamo externo) y el Fondo Europeo de Desarrollo (FEO).

El NDICI descansa en tres pilares: geográfico, temático y de respuesta rápida. Dentro del pilar geográfico, la partida destinada a la Vecindad Oriental y Meridional (1 6 países) es la segunda más importante, detrás de África subsahariana, estando dotada con unos 19.323 millones de euros.

En el marco de la PEV también existen instrumentos que buscan fomentar el acceso a los mercados (apuesta por la negociación de acuerdos de libre comercio de nueva generación), así como instrumentos para afrontar otro de los grandes retos en este contexto: la movilidad de las personas y la gestión de la migración. Sobre esto último, se han celebrado asociaciones de movilidad y mecanismos de· facilitación o liberalización de visados con algunos de los países. Valoración

Pese a la relevancia en términos económicos y geopolíticos de esta política, casi veinte años después de su puesta en marcha ésta no ha logrado aún su objetivo fundamental, crear una estabilidad duradera en sus estados vecinos a través del incremento de las interdependencias económicas y sociales.

Esto se debe a la elevada complejidad del objetivo último de esta política. Por un lado, los países a los que se enfoca se enfrentan a un constante dilema, debido a su situación geoestratégica, entre aproximarse y formar parte del proyecto de sus vecinos europeos occidentales o de sus vecinos asiáticos. Por otro lado, dado que la UE no cuenta con competencias exclusivas en el marco de la PESC se encuentra con el reto de conciliar los intereses, en algunas ocasiones enfrentados, de sus estados miembros, dificultando la formulación de una estrategia común hacia sus países vecinos. Adicionalmente es una política que cuenta con unos recursos posiblemente insuficientes dada la ambición del objetivo, y que en ocasiones es rechazada por algunos de los países a los que se enfoca al considerarse una especie de "imperialismo suave" que busca transmitir los valores  europeos.

\paragraph{Espacio Económico Europeo (EEE, 1992)}
El Acuerdo sobre el Espacio Económico Europeo, firmado en Oporto en mayo de 1992 y en vigor desde el 1 de enero de 1994, extiende a Islandia, Liechtenstein y Noruega (Estados de la Asociación Europea de Libre Comercio, EFTA) las cuatro libertades del mercado interior de la Unión: libre circulación de mercancías, personas, servicios y capitales. Suiza, cuarto miembro de la AELC, rechazó su participación en el EEE mediante referéndum en diciembre de 1992, y mantiene en su lugar una red propia de acuerdos bilaterales sectoriales con la Unión, estructuralmente distinta y más fragmentada.

Sin embargo, no alcanza la Política Agrícola Común ni a la Política Pesquera Común, que quedan fuera del Acuerdo (si bien existen protocolos específicos sobre determinados productos agrícolas transformados y sobre acceso recíproco a caladeros). Tampoco constituye una unión aduanera ya que los Estados del EEE conservan su propia política comercial exterior y sus propios aranceles frente a terceros países, lo que obliga a mantener normas de origen y controles aduaneros en las fronteras entre el EEE y la Unión.\footnote{%
    Esta es la diferencia estructural más relevante frente a la unión aduanera con Turquía: el EEE integra el comercio de mercancías, servicios, personas y capitales, pero no la política comercial exterior; la unión aduanera con Turquía hace exactamente lo contrario.
}

El funcionamiento del EEE descansa sobre el llamado \textit{principio de homogeneidad}: el acervo de la Unión relevante para el mercado interior se incorpora de manera dinámica al ordenamiento de los Estados del EEE a través de decisiones del Comité Mixto del EEE, sin que estos Estados participen en el proceso legislativo de la Unión que da origen a esas normas. 

Esta asimetría (los Estados del EEE deben aplicar una normativa en cuya elaboración no han votado) ha sido descrita en la doctrina como el precio de la integración sin membresía, y es objeto de crítica recurrente por parte de sectores euroescépticos noruegos e islandeses, que la han calificado gráficamente de \textit{democracia por fax}. Para preservar la autonomía del ordenamiento jurídico de la Unión, el EEE se articula sobre una estructura institucional de dos pilares: la vigilancia del cumplimiento del Acuerdo por parte de los Estados del EEE corresponde no a la Comisión Europea, sino al Órgano de Vigilancia de la AELC (ESA, por sus siglas en inglés), y el control judicial no corresponde al TJUE, sino al Tribunal de la AELC, que sigue no obstante muy de cerca la jurisprudencia del TJUE para preservar la homogeneidad de interpretación entre ambos pilares. 

Como contrapartida financiera a los beneficios de acceso al mercado interior, Islandia, Liechtenstein y Noruega contribuyen a la cohesión económica y social de la Unión a través de los mecanismos financieros del EEE y de Noruega (\textit{EEA and Norway Grants}).

\paragraph{Acuerdos de Estabilización y Asociación (AEA): Candidatos a la Adhesión}
El Proceso de Estabilización y Asociación, lanzado a partir del compromiso político asumido en la Cumbre de Tesalónica de 2003, se instrumenta jurídicamente a través de Acuerdos de Estabilización y Asociación (AEA) bilaterales con cada uno de los socios de la región. A fecha de hoy tienen un AEA en vigor Albania, Bosnia y Herzegovina, Kosovo, Macedonia del Norte, Montenegro y Serbia; el de Kosovo, firmado en 2015 y en vigor desde abril de 2016, fue el primer acuerdo internacional suscrito por la Unión tras la entrada en vigor del Tratado de Lisboa, que le confirió personalidad jurídica propia.\footnote{%
    Esto es particularmente interesante desde una perspectiva política: La Unión firmó este último Tratado en su propio nombre, pues hay países de la Unión (como España) que no reconocen su condición de Estado soberano.
}

A diferencia del EEE y de la unión aduanera turca, que son regímenes concebidos como estables en el tiempo, los AEA son por definición transitorios: su función es preparar progresivamente al socio para la adhesión, y no ofrecerle una alternativa permanente a esta. Por tanto, su contenido habitual combina tres elementos: el establecimiento gradual de una zona de libre comercio bilateral con la Unión, el diálogo político estructurado, y, sobre todo, un compromiso de aproximación progresiva de la legislación nacional al acervo comunitario, sujeto a seguimiento y condicionalidad por parte de la Comisión a través de los sucesivos informes anuales de progreso.

Ahora bien, no todos los socios del Proceso de Estabilización y Asociación tienen la misma condición formal ni el mismo grado de avance. Turquía (2004), Macedonia del Norte (2005), Montenegro (2010), Serbia (2012), Albania (2014) y, más recientemente, Bosnia y Herzegovina (2022) ostentan el estatuto de país candidato oficial; Kosovo, en cambio, permanece como candidato potencial únicamente, al no estar su independencia reconocida por cinco Estados miembros de la Unión (España, Chipre, Grecia, Rumanía y Eslovaquia). El grado de avance en la negociación de capítulos del acervo es también muy dispar: Montenegro encabeza el proceso, habiendo cerrado varios capítulos de negociación en 2025 y con el objetivo declarado de concluir las negociaciones de adhesión a finales de 2026, mientras que Serbia no ha abierto ningún capítulo nuevo en los últimos cuatro años, y Bosnia y Herzegovina arrastra, además, obstáculos de naturaleza constitucional derivados del reparto étnico de cargos públicos declarado contrario al Convenio Europeo de Derechos Humanos por el TEDH en el asunto \textit{Sejdić y Finci c. Bosnia y Herzegovina} (2009).

No obstante, la invasión rusa de Ucrania en febrero de 2022 transformó la política de ampliación hacia los Balcanes Occidentales, que pasó de percibirse como un expediente técnico de bajo interés político a considerarse una cuestión de seguridad estratégica, ante el riesgo de que el vacío de influencia europea en la región fuera ocupado por Rusia o por China (que han intensificado sustancialmente su presencia inversora, particularmente en Serbia). En este contexto, la Comisión propuso en 2023 un Plan de Crecimiento para los Balcanes Occidentales que permite a los socios acceder anticipadamente a determinados beneficios del mercado único y a financiación adicional antes de la adhesión formal, condicionado al cumplimiento de reformas, y ha intensificado la celebración de cumbres UE-Balcanes Occidentales, la más reciente, en Tivat (Montenegro), en junio de 2026.\footnote{%
    Consejo Europeo e Infobae, cobertura de la Cumbre UE-Balcanes Occidentales, Tivat, junio de 2026.
}

\paragraph{Acuerdo de Comercio y Cooperación con el Reino Unido (TCA, 2021)}
El Acuerdo de Comercio y Cooperación (TCA) con el Reino Unido, aplicado provisionalmente desde el 1 de enero de 2021 y en vigor con carácter definitivo desde el 1 de mayo de 2021, ocupa una posición singular dentro de esta categoría: no responde a una lógica de integración creciente entre dos economías que se aproximan, como el resto de acuerdos de Nueva Generación, sino a la ordenación de la salida de un antiguo Estado miembro del mercado interior.

El Acuerdo garantiza la ausencia de aranceles y de contingentes para el comercio de mercancías originarias de las partes, pero sujeta ese acceso preferencial al cumplimiento de estrictas normas de origen y, sobre todo, a un conjunto de cláusulas de \textit{level playing field} (no regresión en materia social, laboral, medioambiental y de ayudas de Estado) reforzadas por un mecanismo de reequilibrio (\textit{rebalancing mechanism}) que permite a cualquiera de las partes adoptar medidas arancelarias correctoras, unilateralmente, si constata divergencias regulatorias significativas de la otra parte que afecten sustancialmente al comercio o la inversión bilaterales. 

Este mecanismo, sin precedente en los acuerdos de Nueva Generación anteriores, es una de las singularidades más preguntadas del TCA. La gobernanza del acuerdo corresponde a un Consejo de Asociación y varios comités especializados, y su solución de controversias se articula mediante paneles arbitrales ad hoc.





\subsubsection{Foros y Conferencias Internacionales: SEAE, Representación exterior de la Unión}
Además del esfuerzo de la UE en afianzar un territorio propio y con fronteras estables, la UE quiere jugar un papel relevante en el mundo. La fonna más directa de verlo es mediante el análisis de su participación en foros internacionales. Sin embargo, el papel como actor global de la UE también viene determinado por cuestiones como la fuerza del euro y las acciones de política exterior.

II.II.J. Presencia en foros internacionales 
La presencia de la UE en cada organismo varía en función del diseño institucional de la propia organización y oscila desde la representación única y total de los EEMM (caso de la OMC) a la ausencia y voz a través de los EEMM, que típicamente coordinan sus posiciones (caso del FMI o del Banco Mundial -BM-). Podemos destacar su participación en:

• G20 y G7: Se trata de los principales foros informales de coordinación económica y persiguen el acercamiento de posiciones en materia de política económica para evitar situaciones subóptimas. Sin embargo, mientras que el 020 es un foro con presencia de economías industrializadas y emergentes (y por tanto un l ugar de debate entre posiciones e intereses opuestos), el G7 reagrupa a las principales potencias globales y tienen como finalidad la fijación de posturas comunes. La UE es m iembro de pleno derecho en ambos, junto con tres de sus Estados m iembros, Francia, Alemania e Italia (España es invitado permanente en el 020) y participa a través de la presidenta de la Comisión Europea, y el presidente del Consejo Europeo.

Durante la última reunión del 020, que tuvo lugar en Bali en noviembre de 2022, la UE ha defendido el apoyo a Ucrania, y su compromiso para garantizar la seguridad alimentaria manteniendo el funcionamiento de las cadenas de suministro alimentario, y el acceso a un suministro de energía asequible. Todos los miembros se han reafirmado en su compromisode hacer frente al cambio climático y en la importancia del desarrol lo de las capacidades digitales. Además, han propuesto la creación de un nuevo fondo de intermediación financiera para la prevención, preparación y respuesta frente a pandemias (el «fondo para pandemias»),albergado por el Banco Mundial.
• FMI y BM: La UE no está representada como tal, por lo que su influencia se basa en los derechos de voto de los EEMM.

En el FMI, aunque no cuenta con silla propia la Comisión participa en las reuniones de otoño y primavera de la Junta de Gobernadores, y es observadora del Comité Financiero y Monetario Internacional, los EEMM de la UE con Directores Ejecutivos coordinan sus decisiones en ECOFíN, y al BCE se le invita a asistir a las reuniones del Directorio cuando se debaten temas que conciernen a los países de la moneda única.

En el BM, la Comisión coordina las relaciones de los EEMM con la institución para dar un enfoque uniforme y es observadora.en el Comité para el Desarrollo. Además de en el BM, los EEMM de la UE tienden a fijar posiciones comunes en los bancos regionales de desarrollo. Caso especial es el Banco Europeo de Reconstrucción y Desarrollo (EBRD), donde los países de la UE cuentan con la mayoría del capital.

• BIS y FSB: En cuanto al Banco Internacional de Pagos (BIS), el principal centro de cooperación internacional de bancos centrales, el BCE participa. Y en el Financia! Stability Board (FSB), participan tanto la UE directamente (Comisión y BCE) como los EEMM con presencia en el 020; a través de este foro la UE influyó en el diseño de la arquitectura financiera internacional posterior a la crisis financiera de 2007-2008.

• OCDE: En la OCDE la UE tiene carácter de "cuasi-miembro", ya que participa en el conjunto de trabajos y estrategias, manteniendo una delegación permanente ante esta institución, pero no tiene derecho a voto en las decisiones del Consejo, y no está obligada a las contribuciones presupuestarias.

• ONU: En la ONU, tanto en la Asambl��a como en el Consejo de Seguridad la UE sólo es observador, un rango inferior a los EEMM que están representados en ella (si bien sólo Francia es miembro permanente del Consejo de Seguridad). No obstante, la coordinació entre EEMM y la UE es crucial para defender los intereses comunitarios, aunque en política exterior pueden ser muy diferentes, y la capacidad de decisión es de los Estados. Pero en otros órganos, como la conferencia de NNUU sobre Comercio y Desarrollo. o el Programa de NNUU para el Desarrollo su papel es mayor, al ser materia de su competencia el comercio, la agricultura, la pesca, y el desarrollo.

• OMC: A diferencia de lo que sucede en casi todas las instituciones que hemos visto hasta el momento, en la OMC la UE tiene todo el protagonismo en detrimento de los EEMM. Así, pese a que todos ellos disponen de representantes nacionales en la OMC, las negociaciones en las Rondas las lleva a cabo la Comisión ya que, con carácter general, la política comercial es competencia exclusiva de la UE. Actualmente, ante el estancamiento del proceso multilateral, la UE aspira a ser un mediador central, un constructor de consensos.



\subsubsection{Política Común de Defensa: Autonomía estratégica, industria de defensa}
Curiosamente, las cuestiones de seguridad y defensa han sido siempre objeto de una especial atención en los Tratados, si bien es cierto que no siempre para incluirlas, sino para hacerlas objeto de un tratamiento singularizado en relación con otros aspectos materiales de la PESC. 

Aunque este concepto no tiene una definición única, "autonomía estratégica" se refiere a la capacidad de la UE de actuar de manera autónoma, al menos en productos o servicios considerados críticos, o coordinada con terceros si fuera conveniente.\footnote{%
    Este término se usó por primera vez en las Conclusiones del Consejo Europeo sobre la política común de seguridad y defensa de diciembre de 2013. Francia lo asocia a soberanía. frente a países con mayor tradición  trasatlántica como Países Bajos, o fronterizos con Rusia como Finlandia o Estonia. de tradición más liberal.
} Según la Comisión (2021), sería la capacidad de la UE de tomar sus propias decisiones y dar forma al mundo que la rodea a través del liderazgo y el compromiso, reflejando sus intereses y valores estratégicos, a largo plazo".\footnote{Trade policy review: an open. sustained and assertive trade policy
}
 
En un contexto internacional crecientemente marcado por diferentes tensiones geopolíticas y cambios económicos y tecnológicos que dificultan la cooperación internacional (ascenso de China como potencia, impacto de la pandcmia del Covid-19, rupturas en las cadenas globales de suministro de bienes y servicios básicos, dependencia de materias primas y falta de reservas estratégicas, desequilibrios globales en los mercados energéticos, inflación y amenaza a la seguridad derivada de la invasión rusa de Ucrania), conseguir mantener la autonomía estratégica al tiempo que se preserva una economía abierta es un objetivo clave de la UE.\footnote{%
    En respuesta a las voces que alertan de que la autonomía estratégica podría estar fomentando cierto proteccionismo y poniendo en riesgo la competencia en el mercado interior.
} Este nuevo contexto pone de relieve la pérdida de relevancia del multilateralismo en general y de la UE en particular, que tradicionalmente ha ejercido un poder blando, no coactivo y civil, basado en la cooperación con otros países y en el establecimiento de estándares europeos por parte de terceros países (el llamado "Efecto Bruselas").

La autonomía estratégica cuyo origen se sitúa en el ámbito de la seguridad/defensa, se ha ido ampliando posteriormente y actualmente también abarca las esferas económica, sanitaria, industrial y tecnológica, comercial y energética, ya que en el fondo todas ellas están interrelacionadas.

En el ámbito económico, la Declaración de Versalles de l 0-11 de marzo de 2022 señala como prioridades esenciales de la UE:
• Aumentar la inversión en capacidades de defensa en los Estados miembros, e impulsar la colaboración a nivel UE, incluyendo las compras públicas conjuntas
• Impulsar la independencia energética. Se está poniendo en marcha la eliminación gradual de las importaciones rusas de petróleo, gas y carbón. Además, el nuevo Reglamento RepowerEU, que dota de financiación adicional a los planes nacionales de recuperación y resiliencia servirá para ahorrar energía, diversificar las importaciones energéticas y acelerar la sustitución de combustibles fósiles por energías renovables.
• Dotar a la UE de una base económica robusta que respalde al euro frente al dólar como moneda de reserva internacional.

Éstas se vienen a sumar a las planteadas a su vez cuando se aprobó el instrumento NextGeneration EU: acelerar las tran siciones gemelas, esto es, la transformación energética de cara a cumplir los objetivos climáticos de aquí a 2050 y la digitalización de la economía. Todo ello requiere aunar voluntades entre los Estados miembros para definir intereses comunes y crear instrumentos eficaces.

En las conclusiones del Consejo de la UE de marzo de 2022 destacan entre otros, dos objetivos más concretos:\footnote{%
    Council Conclusions on the EU's economic and financia! strategic autonomy: one year after theCommission's Communication
} i) asegurar un papel internacional del euro más fuerte profundizando la Unión Económica y Monetaria y promoviendo el uso de activos denominados en euros; y ii) asegurar la resiliencia del sector financiero, completando la Unión Bancaria y desarrol lando la Unión de Mercados de Capitales.

El mercado de defensa de la UE se encuentra aún fragmentado: los sistemas de defensa nacionales presentan escasos avances en materia de interoperabilidad y pocos compromisos de invertir el nuevo gasto en defensa de forma conjunta. Se observa una importante duplicidad de esfuerzos (presupuestarios, de recursos humanos e infraestructuras). Ello puede explicarse en parte porque estos mercados han estado tradicionalmente protegidos de la competencia externa por diversos motivos: 1) evitar la dependencia de terceros países y 2 ) desarrollar capacidades propias a medio y largo plazo para crear nuevas tecnologías, puestos de trabajo y evitar problemas de seguridad en el suministro de piezas, repuestos y en el mantenimiento de los sistemas. Además, las distintas percepciones de seguridad, culturas estratégicas y restricciones ·burocráticas de los países de la UE, junto a la secular escasez presupuestaria, son !imitadores evidentes del desarrollo industrial, estratégico y militar.

La Política Común de Seguridad y Defensa (PCSD) forma parte de la PESC. La nueva PCSD busca fomentar la cooperación intraeuropea entre empresas para mejorar la eficiencia e impulsar las sinergias existentes.

Según la encuesta Eurobarómetro especial sobre seguridad y defensa de 2017, el 55\% de la población encuestada estuvo a favor de la creación de un ejército de la UE. A van zar hacia una unión de seguridad y defensa es una de las prioridades de la Comisión actual. No obstante. persisten algunos de los problemas subyacentes a la creación de un ejército europeo como las diferentes valoraciones que los Estados miembros hacen de la urgencia de responder a una determinada crisis, la financiación de las operaciones, o el diseño actual de la PESC, cuya lógica es intergubernamental (por lo que se requiere de unanimidad en las votaciones) y que sólo financia misiones de carácter civil.\footnote{%
    Por ejemplo, Hungría ha ejercido su veto para la adopción de algunos paquetes de sanciones frente a Rusia.
} 

Aunque aún no existe un "ejército de la UE" y la defensa sigue siendo una prerrogativa nacional exclusiva, lo que sin duda explica las duplicidades existentes, la UE ha dado grandes pasos para impulsar la cooperación. La Brújula Estratégica avalada por el Consejo Europeo en marzo de 20 22 es la nueva estrategia de seguridad y defensa con horizonte 2030: puede asimilarse a un ·'libro blanco'· de la defensa. Partiendo de un análisis del entorno que se ha vuelto más hostil en los últimos años (aumento de ciberataques, campañas de desinformación, y otras amenazas híbridas) propone por primera vez una guía para la acción colectiva de los EEMM.\footnote{%
    Cuenta con propuestas en cuatro ámbitos: i) actuar con mayor rapidez y decisión ante las crisis: ii) proteger a los ciudadanos frente a las amenazas cambiantes: iii) invertir en las capacidades y tecnologías que necesitamos; iv) asociarse con otros para alcanzar objetivos comunes. En total incluye más de 80acciones concretas.
} Entre sus principales elementos. destaca la nueva Capacidad/Fuerza de Despliegue Rápido de la UE (EU RDC) que consistirá en agrupaciones tácticas sustancialmente modificadas y fuerzas y capacidades militares, para disponer de una fuerza modular de hasta 5.000 efectivos. La VE realizará las primeras maniobras militares en el segundo semestre de 20 23 y se espera que esté plenamente operativo en 20 25.

El tercer objetivo de la Brújula Estratégica pone en valor la importancia de la industria militar europea. Los EEMM se han comprometido a incrementar sustancialmente su gasto en defensa (unos 70.000 millones de euros más para 20 25) para reforzar las capacidades militares y civiles (Cooperación Estructurada Permanente) y la base tecnológica e industrial de la defensa europea (EDTIB, por sus siglas en inglés).\footnote{%
    En la Cumbre de Gales de la OTAN de 2014, los miembros se comprometieron a aumentar el gasto en defensa hasta el 2\% del PIB para 2024. Según el último informe anual del Secretario General de la OTAN de 2022. solo 7 de los 30 países lo cumplen (Grecia Estados Unidos, Lituania, Polonia Reino Unido, Estonia y Letonia). En España representa el 1 .09\%.

Francia y Alemania anunciaron en 2022 aumentos de inversión sustanciales en este sector a corto y medio plazo. La UE por su parte acordó financiar con 1.200 millones de euros 61 proyectos para apoyar la capacidad de defensa y la tecnología crítica.
}

Las Conclusiones del Consejo Europeo de diciembre de 20 22 rei!eran que la UE está asumiendo mayor responsabilidad por su propia seguridad, siguiendo la "Brújula Estratégica" que contribuirá al refuerzo de la base industrial y tecnológica del sector europeo de la defensa. Tras hacer balance de los progresos, pide a las instituciones europeas principalmente: i) que aceleren la aprobación de la legislación sobre contratación pública común en materia de defensa (EDIRP A por sus siglas en inglés), en particular a la luz del apoyo prestado a Ucrania;\footnote{%
    La Comisión propuso enjulio de 2022 dotar EDIRPA con 500 millones de euros. Este nuevo instrumento de compras públicas conjuntas complementará al Fondo Europeo de Defensa dotado con 8.000 millones de euros para el periodo 2021 -2027 y al Fondo Europeo de Apoyo a la Paz dotado con 2.000 millones.
} ii) que se presente una propuesta de Programa Europeo de Inversiones en Defensa para reforzar la capacidad y la eficacia de la defensa europea, así como la resistencia del sector tecnológico, incluidas las pyme, iii) acelerar la ejecución de los proyectos de infraestructuras de movilidad militar (incluidas las de doble uso), iv) mayor inversión en ciberseguridad, conectividad, y reforzar las infraestructuras criticas y v) la rápida aplicación de la caja de herramientas híbrida de la UE.

La reacción de la UE a la invasión rusa de Ucrania se ha centrado en diversos paquetes de sanciones (el 10° paquete fue aprobado en febrero de 20 23) y en el apoyo financiero y militar a Ucrania.\footnote{%
    Según el informe anual de progreso de la Brújula Estratégica de marzo 2023. la UE ha proporcionado una ayuda global colectiva a Ucrania, incluida la ayuda macro1inanciera y humanitaria de al menos 67.000 millones de euros.
} Esta guerra ha transformado el enfoque tradicional de la PCSD, más centrado en priorizar las capacidades y en la gestión de crisis a través de operaciones de intensidad militar baja o media fuera del territorio europeo, y ha evidenciado la difícil adaptación de la industria de defensa, que funciona bajo pedidos, cuando se produce un aumento súbito de la demanda. Ello evidenciaría la necesidad de potenciar la cooperación intra-UE en este campo.







\begin{specialblock}{Papel Internacional del Euro: Resilencia del Sector Financiero}
    En 2019 se conmemoró el 20° aniversario del euro y en 2023 tras incorporarse Croacia, la eurozona cuenta con veinte EEMM y más de 340 millones de personas. Según los resultados de la encuesta trienal de 2022 del BIS '1, el euro se mantiene por detrás del dólar como segunda moneda más utilizada en los pagos internacionales y como moneda de reserva (30,5\% de todas las operaciones, ligeramente inferior al 32\% de 20 19 frente al 88\% de todas las Qperaciones en abril de 2022, respectivamente). Además, otros sesenta países y territorios, en los que viven unos 175 millones de personas, han decidido utilizar el euro como moneda o vincular su propia moneda a él.\footnote{
    Los microestados República de San Marino, Ciudad del Vaticano. Principado de Mónaco y Andorra utilizan el euro como moneda ofü:ial. Los países cuyas monedas están estrictamente vinculadas al euro son: Dinamarca, Bu lgaria. Serbia, Bosnia-Herzeg ovina. Macedonia, Rumanía. Albania. Túnez y Marruecos.
    }
        
    El uso masivo de una moneda genera influencia por parte de su emisor. Por ello, reforzar su papel internacional contribuirá a forta lecer el liderazgo europeo mundial y generará múltiples beneficios que la Comisión ya detectó en la Comunicación publicada en 2018:
    • Menor coste y riesgo del comercio internacional para las empresas europeas
    (especialmente las pymes)
    • Más posibilidades de elección para los operadores de mercado de todo el mundo.
    • Tipos de interés más baj os para los hogares, las empresas y los Estados miembros europeos. Un euro más atractivo como depósito seguro de valor reduce el tipo de interés (o rendimiento) exigido por los inversores.
    • Acceso más fiable a la financiación para las empresas y los gobiernos europeos, incluso en periodos de inestabilidad financiera externa, ya que los mercados financieros europeos podrían hacerse más profundos, líquidos e integrados.
    • Mayor autonomía de los consumidores y las empresas europeas, permitiéndoles pagar o recibir pagos por su comercio internacional y financiarse con una menor exposición a las acciones legales emprendidas por jurisdicciones de terceros países, como las sanciones extraterritoriales.
    • Mejora de la resiliencia del sistema financiero y la economía internacionales, haciéndolos menos vulnerables a las perturbaciones de los tipos de cambio. En enero de 2021 la Comisión reforzó su interés en utilizar el euro como una herram ienta de influencia económica mundial, y publicó una nueva Comunicación sobre el lanzamiento de una nueva estrategia para estimular la apertura, la solidez y la resiliencia del sistema económico y financiero de la UE.\footnote{%
        En el Anexo 4 se resumen las distintas iniciativas para potenciar el papel internacional del euro.
    } Se prevén medidas en el ámbito financiero (intra-UE e internacional) y energético:
    • Apoyar el desarrollo de instrumentos y referencias denominados en euros y fomentar su condición de moneda de referencia internacional en los sectores de la energía y las materias primas, incluidos los vectores energéticos incipientes como el hidrógeno.
    • La emisión de bonos de alta calidad denominados en euros en el marco de la NextGeneration EU añadirá una profundidad y liquidez significativas a los mercados de capitales de la UE en los próximos años y los hará, a ellos y al euro, más atractivos para los inversores.
    • Promover las finanzas sostenibles es también una oportunidad para convertir los mercados financieros de la UE en un centro mundial de "finanzas verdes", reforzando el euro como moneda por defecto para los productos financieros sostenibles.\footnote{%
        La propia Comisión financiará hasta 250.000 M€ (o el 30\%) de Next-Generation EU mediante la emisión de bonos verdes Next-Generation EU. Esto la convertirá en el mayor emisor de bonos verdes del mundo.
    }
    • La Comisión también buscará posibil idades de ampliar el papel del Régimen Comunitario de Comercio de Derechos de Emisión (ETS por sus siglas en inglés) para maximizar sus resultados medioambientales y apoyar la actividad comercial del ETS en la UE. Asimismo, la Comisión presentó en diciembre de 2021 una propuesta para establecer los ingresos derivados del ETS como recurso propio de la UE.
    • Además de todo esto, la Comisión también seguirá apoyando el trabajo del BCE sobre una posible introducción de un euro digital, como complemento del efectivo. Asimismo, no debe olvidarse el papel protagonista del BCE para apoyar la recuperación económica global tras la pandemia de Covid- 19, a través de su participación en foros internacionales para coordinar la respuesta a la crisis de manera multilateral, o de la apertura de líneas de swap de divisa con otros bancos centrales para la correcta provisión de euros a los mercados internacionales

    Por otra parte, es importante señalar que la UE está actuando para reducir la dependencia financiera del exterior que se observa en los siguientes aspectos: 1) naturaleza estratégica de algunas de las adquisiciones chinas en Europa y la opacidad de las prácticas de inversión chinas en el extranjero, 2) dependencia excesiva de los participantes en los mercados financieros de la UE de servicios de pagos, compensación y liquidación u otros servicios críticos para las infraestructuras de mercado ofrecidos por proveedores de terceros países, con los consiguientes riesgos operacionales y para la estabilidad financiera de la zona euro,. 3) ascenso de grandes empresas tecnológicas de térceros países en el mercado de servicios financieros de la UE y la creciente interconexión entre este mercado y los proveedores de servicios de criptomonedas (muchos de ellos residentes fuera de la UE), y 4) el elevado grado de apertura. de la UE, junto con su papel relativamente reducido en las redes financieras y de pagos internacionales, hace a la Unión particularmente vulnerable a sanciones secundarias impuestas por los países centrales en dichas redes (extraterritorialidad de las sanciones de Estados Unidos). En concreto, la UE está avanzando en la identificación de esas debilidades a través del informe sobre interdependencias estratégicas, y ha puesto en marcha iniciativas más amplias como el plan de acción sobre materias primas críticas, la iniciativa RepowerEU, y los planes de digitalización. Asimismo, la lJE monitoriza los flujos de inversión directa extranjera (ver infra) y está trabajando para reforzar las infraestructuras financieras. Por último, deben mencionarse las medidas destinadas a preservar el buen funcionamiento del mercado interior, como la regulación sobre subvenciones extranjeras o el Mecanismo de Ajuste en frontera del Carbono.
\end{specialblock}



\subsubsection{Medidas restrictivas}
Otro capítulo que supone un doble anclaje en las disposiciones PESC y la regulación de las otras competencias de la Unión son las medidas restrictivas. La decisión de adoptar medidas que puedan suponer la interrupción o la reducción, total o parcial, de las relaciones económicas y financieras con uno o varios terceros países que, eventualmente, siguiendo las nuevas tendencias, pueden ser medidas restrictivas contra personas físicas o jurídicas, grupos o entidades no estatales, se adoptan en el ámbito PESC por unanimidad, pero su ejecución se lleva a cabo por el Consejo por mayoría cualificada, según se establece en el artículo 215 TFUE en el ámbito de la acción exterior. En este último caso, según dispone el artículo 275 TFUE, el TJ tiene excepcionalmente competencia sobre la propia decisión PESC, ya que conoce sobre los recursos interpuestos en las condiciones contempladas en el artículo 263 TFUE (relativo al control de la legalidad de las decisiones adoptadas por el Consejo en virtud del Cap. 2 del Título V del TUE). Es una novedad del Tratado de Lisboa que tiene su origen en el desolador panorama en materia de garantías que venía causando en este contexto la situación hasta ahora vigente y que se plasmó en una confusa, artificiosa y decepcionante saga de sentencias del TG y del TJ16. No parece que las normas adoptadas aquí y en el ELSJ sobre restricciones en capitales y pagos (art. 75 TFUE) ofrezcandemasiada tranquilidad sobre la solución definitiva del problema17.

Desde 2014, la UE viene imponiendo progresivamente medidas restrictivas contra Rusia en respuesta a la anexión ilegal de Crimea en 2014, el ataque militar sin precedentes y no provocado de Rusia contra Ucrania en 2022 y la anexión ilegal de 4 regiones ucranianas en 2022. En este contexto la UE también ha adoptada medidas restrictivas contra Bielorrusia por su participación en la invasión de Ucrania e Irán por el uso de drones iraníes en esta situación. Estas medidas restrictivas pueden dividirse en individuales, sanciones económicas, restricciones a los medios de comunicación, medidas diplomáticas, medidas en materia de visados, restricciones a las relaciones económicas y medidas relativas a la cooperación económica. Las medidas restrictivas individuales hacen referencia principalmente a la inmovilización de bienes y restricciones de viaje por actos que menoscaban la integridad territorial, la soberanía y la independencia de Ucrania.

La lista de personas y entidades sancionadas es revisada y renovada periódicamente e incluye, entre otros, a Vladimir Putin, miembros de la Duma Estatal rusa, militares y altos funcionarios, bancos y entidades financieras, partidos políticos, etc. Estas medidas se prorrogaron por última vez hasta el 1 5 de marzo de 2023 (adoptadas por primera vez en 2014). Las sanciones económicas también vienen imponiéndose desde 2014 y se aplican a los intercambios con Rusia de sectores económicos determinados. En 2015 la UE decidió la vinculación de las sanciones a la plena aplicación de los Acuerdos de Minsk.\footnote{%
    Los acuerdos de Minsk son dos pactos firmados en 20 1 4 y 20 1 5 por Rusia Ucrania. la República Popular de Donetsk y la República Popular de Lugansk para poner tin a la guerra del Donbás (Ucrania). Éstos se han ido incumpliendo por ambas partes.
} Esto ha llevado a que las sanciones económicas se hayan ido prorrogando cada 6 meses desde 2016, siempre tras • evaluar la aplicación de dichos acuerdos.

Estas sanciones van dirigidas a los sectores comercial, financiero, energético, de transporte, de defensa y tecnológico.
Sector financiero: un acceso limitado a los mercados de capitales primario y secundario de la UE para determinados bancos y empresas rusas, la prohibición de las transacciones con el Banco Central de Rusia y el Banco Central de Bielorrusia, y con el Banco de Desarrollo Regional ruso, la prohibición de acceso al sistema SWIFT para determinados bancos rusos y bielorrusos, la prohibición del suministro de billetes denominados en euros a Rusia y Bielorrusia, la prohibición de la financiación o la inversión públicas en Rusia, la prohibición de invertir en proyectos cofinanciados por el Fondo Ruso de Inversión Directa y de contribuir a estos.
Energía: la prohibición de las importaciones de carbón procedentes de Rusia, la prohibición de las importaciones de petróleo procedentes de Rusia, con excepciones limitadas, la limitación de precios en relación con el transporte marítimo de petróleo, establecido en 60 USO por barril, la prohibición de exportación a Rusia de bienes y tecnologías del sector del refino de petróleo, la prohibición de nuevas inversiones en el sector ruso de la energía y la minería.
Transporte: el cierre del espacio aéreo de la UE a todas las aeronaves rusas o matriculadas en Rusia, la prohibición del transporte marítimo de petróleo ruso a terceros países (si se compra por encima del límite de precio), el cierre de los puertos de la UE a los buques rusos, la prohibición de entrar en la UE a los transportistas por carretera rusos y bielorrusos, la prohibición de la exportación a Rusia de bienes y tecnología de las industrias de la aviación, marítima y espacial.
Defensa: la prohibición de la exportación a Rusia de productos de doble uso y tecnología, la prohibición de la exportación a Rusia de motores de drones, la prohibición del comercio de armas y de armas de fuego de uso civil, la prohibición del comercio de municiones, vehículos militares y equipo paramilitar.

Materias primas y otros bienes: la prohibición de las importaciones de hierro, acero, madera, cemento y plásticos procedentes de Rusia, la prohibición de la importación de oro procedente de Rusia, la prohibición de las importaciones de pescado y marisco, bebidas alcohólicas, cigarrillos, joyería y cosméticos procedentes de Rusia.

Adicionalmente, en 2022 la UE suspendió las licencias y actividades de radiodifusión de múltiples emisoras rusas, como Sputnik, debido a su utilización por parte del Gobierno ruso como instrumentos para manipular la información y promover la desinformación con el objetivo de desestabilizar a los países vecinos de Rusia y a la UE.

Las medidas diplomáticas han constado de la suspensión de las cumbres bilaterales periódicas UE- Rusia desde 2014, la suspensión de las conversaciones en materia de visados, la celebración de reuniones del G7 como sustitución del G8 (que incluía a Rusia) y la suspensión de las negociaciones de adhesión de Rusia a la OCDE y la AIE (Agencia Internacional de la Enrgía). En materia de visados en 2022 la UE estableció que los diplomáticos, funcionarios y empresarios rusos dejaban de poder acogerse a las disposiciones sobre facilitación de visados que permiten un acceso privilegiado a la UE, suspendiendo totalmente en septiembre el Acuerdo de facilitación de visados UE-Rusia.

La UE también ha adoptado medidas en respuesta a las injerencias de Rusia en distintas regiones rusas (Crimea, Sebastopol, Donetsk, Luhansk, Zaporitia y Jersón) que aplican a los nacionales y empresas con sede en la UE y que incluyen la prohibición de importación de bienes, restricciones · al comercio y a la inversiones en determinados sectores económicos y proyectos de infraestructuras, la prohibición de prestar servicios turísticos y la prohición de exportación de determinados bienes y tecnologías.

Por último, respecto a las medidas relacionadas con la cooperación económica, desde 20 1 4, el Banco Europeo de Inversiones suspendió la firma de nuevas operaciones de financiación en Rusia, los EEMM acordaron coordinar sus posiciones en el Banco Europeo de Reconstrucción y Desarrollo (BERD) con el objetivo de suspender la financiación de nuevas operaciones.































\clearpage
\section{Acción Exterior de la Unión: Cooperación al desarrollo}
\puntosclave{
    La Agenda 2030 para el Desarrollo Sostenible centrada en las 5P (personas, planeta, prosperidad, paz y asociación) actúa como una brújula que guía las medidas adoptadas por la UE para alcanzar los 17 ODS. 
    
    La coherencia y la eficacia de las distintas acciones que forman parte de la cooperación al Desarrollo son fundamentales para aprovechar sinergias y maximizar el impacto. también juegan un papel crucial la prográmación geográfica y la evaluación de resultados.
    
    El nuevo Instrumento Europa Global dentro del Marco Financiero Plurianual 2021-2027 establece la financiación concreta en función de las prioridades geográficas, temáticas y de respuesta rápida. 
    
    Las Iniciativas Equipo Europa permiten identificar prioridades críticas en las que concentrar los esfuerzos financieros y técnicos para apoyar el Desarrollo de un país en un determinado ámbito. 
    
    La UE se mantiene como líder mundial de AOD destinando la mayor parte de los fondos a África Subsahariana.
}


Desde el Tratado de Maastricht, la UE se ve otorgada una nueva, más amplia y congruente cobertura jurídica a esta actividad, las competencias en materia de cooperación económica, financiera y técnica son caracterizadas como complementarias de las de los Estados miembros y coherentes con las de la cooperación al desarrollo. La Ayuda humanitaria es una competencia complementaria que persigue la prestación de asistencia y socorro por parte de la Unión a las poblaciones de terceros países víctimas de catástrofes naturales o de origen humano. En el ejercicio de esta competencia la Unión se regirá conforme a los principios del Derecho Internacional de imparcialidad, neutralidad y no discriminación;\footnote{%
    La cooperación al desarrollo fue formalizada en el Tratado de Maastricht. Esto no significa que hubiese estado ausente con anterioridad. En realidad, desde hacía tiempo, la actividad de cooperación al desarrollo constituía un componente esencial de las relaciones exteriores de la Unión, que se había ido construyendo progresivamente por distintas vías a partir de las disposiciones relativas a los países y territorios de Ultramar, del mecanismo de la asociación prevista en el hoy artículo 217 TFUE e, incluso, mediante la adopción de medidas presupuestarias, entre otras.
}


\subsection{Cooperación para el Desarrollo: }
Las acciones de la Unión en materia de cooperación al desarrollo se justifican por razones históricas, económicas y políticas distintas, se encauzan a través de medios y mecanismos diferentes y, naturalmente, conducen a niveles o grados de cooperación también distintos.\footnote{%
    El origen de la cooperación al desarrollo se sitúa en el Tratado CEE, cuya finalidad era articular las relaciones con los países vinculados histórica, política y económicamente a los Estados miembros de las CCEE. Pero, superando esa primera finalidad, pronto se convirtieron en el embrión de un sistema comunitario de cooperación internacional para el desarrollo (ABELLÁN HONRUBIA, 1982: 192).
} Por tanto, la instrumentación jurídica en materia de cooperación al desarrollo no es uniforme ni tampoco unitaria. 

En un sentido estricto, la Cooperación para el Desarrollo se enmarca en los arts. $208$ a $211$ TFUE (Quinta Parte, Título III), donde se establecen los objetivos, instrumentos y procedimientos en materia de cooperación al desarrollo así como, naturalmente, el alcance y naturaleza de las competencias de la Unión. Así, se establece como objetivo principal \textit{la reducción y, finalmente, la erradicación de la pobreza}, olvidando pues objetivos como el desarrollo económico y social duradero u otros más políticos. Por otro lado, la UE en esta materia goza de una competencia complementaria con los Estados miembros. De hecho, la calificación de esa política como \textbf{complementaria} condiciona, naturalmente, el alcance y contenido de la competencia de la Unión.\footnote{%
    Se ha interpretado entendiendo la complementariedad como apoyo mutuo, en el sentido de que la política de la UE tiene, al igual que las de los Estados, una entidad y un dinamismo propios que deben favorecer la interacción natural entre ella misma y esas políticas nacionales y cubrir el déficit de articulación existente entre ellas, hasta entonces, con objeto de ofrecer una ayuda al desarrollo global por parte de la Unión y sus Estados miembros (FLAESCH-MOUGIN, 1993: 366).
}

No obstante, siempre deberán observar, como todo, las disposiciones generales que vertebran las acciones de la Unión.\footnote{%
    Esta cláusula se articula a través de dos tipos de medidas: por un lado, acciones positivas dirigidas a fomentar el respeto de los principios democráticos y de los derechos y libertades fundamentales y, por otro lado, acciones negativas como la reducción o supresión de las ayudas en caso de violación de dichos principios o derechos y, en su caso, la suspensión o terminación de los acuerdos con terceros Estados con contenido de cooperación. En ocasiones, los intereses económicos y comerciales prevalecen sobre esas exigencias dando un sesgo indebido a esta política de la UE. Por otra parte, los Estados terceros no aceptaban de buen grado la inclusión de estas cláusulas rechazando de plano la condicionalidad política. Sobre la relevancia de la inserción en el Tratado de esta condicionalidad política se ha pronunciado el TJUE en una importante sentencia de 3 de diciembre de 1996 relativa al Acuerdo de Cooperación con la India.
}

En definitiva, y en linea con lo que decíamos sobre la pluralidad de instrumentos, se ha distinguido tradicionalmente entre los casos en los que la Unión ha recurrido a la adopción: de (i) \textbf{medidas autónomas}; y/o, (ii) conclusión de \textbf{acuerdos internacionales de asociación}. Esta pluralidad de instrumentos se acompaña de una diversificación en el plano de la financiación de las acciones de cooperación al desarrollo, que, según los casos, corren a cargo del presupuesto de la Unión o de contribuciones de los Estados miembros.



\subsubsection{Medidas de carácter autonómo (\textit{ex} art. $209.1$): Programas para el desarrollo}
Las medidas de carácter autónomo (art. $209.1$) se adoptarán siguiendo el procedimiento legislativo ordinario, 

\paragraph{MFP 2021-2027: Programa Europa Global}
Con ocasión del nuevo MFP 2021 -27, la UE decidió reestructura su política de cooperación exterior, con el fin de hacerla más efectiva y explotar mejor sus palancas. Así, el nuevo Instrumento de Vecindad, Desarrollo y Cooperación Internacional (NDICI), denominado "Europa Global" se convierte en el centro del sistema.

El nuevo programa remplaza más de diez instrumentos financieros del periodo presupuestario anterior (2014-2020), entre ellos el Fondo Europeo de Desarrollo (FED) para la cooperación con los países de ACP. Éste último era el más relevante en volumen, pero se financiaba directamente por los EEMM y no por el presupuesto europeo (y por lo tanto con un papel reducido del Parlamento Europeo). A su vez, dichos instrumentos estaban fragmentados por zonas geográficas (el FED para los países ACP; el Instrumento Europeo de Vecindad y Asociación para nuestros vecinos del Sur y el Este; el Instrumento de Ayuda de Preadhesión; el Instrumento de Cooperación al Desarrollo para América Latina, Asia y África del Sur; etc.) y por temas {lnstrumento Europeo para la Democracia y los Derechos Humanos; Instrumento en pro de la Estabilidad y la Paz; Instrumento de Colaboración; etc.).

Con el establecimiento de una ;'Europa Global", la UE busca armonizar y simplificar sus instrumentos de cooperación exterior y AOD, poniendo fin a varias décadas de fragmentación y complejidad. A partir de ahora, el alcance del programa es universal (sus planes con base  geográfica solo excluyen a los países y territorios de preadhesión y de ultramar, que cuentan consu propio instrumento).

Teniendo en cuenta estos elementos, la Comisión tiene la ambición, con la ayuda de Europa Global, de ser más geopolítica para el periodo 2021 -2027, estableciendo asociaciones con todos los países, regiones y actores internacionales y, basándose en su experiencia ante la crisis causada por la pandemia de la Covid-19, trabajar más estrechamente con sus EEMM y otros socios europeos (o que comparten sus mismos objetivos), dentro de un contexto de apoyo a la multil_ateralidad.


Pilares de Europa Global. Son tres los pilares del nuevo sistema: geográfico, temático y de respuesta rápida. 

El pilar geográfico consta de programas geográficos sobre vecindad europea (ya estudiados), África Subsahariana, Asia y el Pacífico y el continente americano y el Caribe. El propio documento estratégico señala el ámbito de aplicación material de los programas geográficos y defme las áreas de cooperación y cómo las acciones se deben llevar a cabo en todas las zonas geográficas, resaltando las peculiaridades para la zona de vecindad34•

Con la adopción en 2018 de sus propuestas en materia exterior para 202 1 -2027, la Comisión declaró que África, la zona de la Vecindad europea y los Balcanes Occidentales (éstos cubiertospor IPA lll) serían sus zonas preferentes para la cooperación, lo que fue confirmado por los acuerdos sobre el nuevo MFP 202 1 -27. Estas prioridades se traducen en una mayor cantidad de recursos financieros para estas zonas geográficas, sin abandonar otras regiones del planeta como Asia y el Pacífico o América y el Caribe.

En términos de dotación, este pilar representa el 75% del total (60.388 millones de euros), distribuidos en las siguientes partidas:
• África Subsahariana 29.18 1 M€
• Asia y Pacífico 8.489M€
• América y Caribe 3.395M€ 

El pilar temático se centra en los desafíos mundiales, principalmente a través de programas temáticos especializados sobre derechos humanos y democracia, organizaciones de la sociedad civil, la estabilidad y la paz, y desafíos mundiales, que abarcan temas como salud, educación y formación, mujeres y niños, trabajo digno y protección social, cultura, migración, medio ambiente y cambio climático, energía sostenible, crecimiento sostenible e integrador, sector privado y autoridades locales. Estos programas son complementarios de los programas geográficos, y tienen cobertura mundial.\footnote{%
    En consonancia. se pueden dedicar ahora sus recursos a aquel las actividades que por su naturaleza son realmente globales y no pueden ser circunscritas a una zona geográfica en concreto, al menos en su diseño inicial, por ejemplo las contribuciones a GA VI o al Fondo Mundial para la lucha contra el SIDA. la Tuberculosis y la Malaria. Por otro lado. el incremento del programa temático para paz. estabilidad y prevención de conllictos subraya de nuevo la importancia de la seguridad en la agenda de la Unión Europea.
} Su dotación es de 6.358 millones de euros (8\% del total). Las áreas prioritarias son de carácter transnacional, y forman el núcleo de los programas con las regiones mencionadas:
• transiciones verde y digital
• crecimiento sostenible y creación de empleo
• gobemanza, paz y seguridad, cultura
• migraciones y desplazamientos forzosos
• salud
• educación, investigación y capacidades

El \textbf{pilar de respuesta rápida} está dedicado a la creación de capacidades de respuesta rápida para la gestión de crisis, la prevención de conflictos y la consolidación de la paz; el refuerzo de laresiliencia y la vinculación entre las medidas humanitarias y de desarrollo; y las prioridades y necesidades en las acciones de política exterior. Este componente también tiene cobertura mundial. Dentro de este componente, no se necesita programación; su aplicación consiste en la adopción directa de medidas de ayuda excepcionales, planes de acción y medidas individuales que se deciden dependiendo de las necesidades y que no pueden predecirse a largo plazo. Su dotación es de 3 . 1 8 2 millones de euros.

Finalmente, además de los tres pilares, se crea la reserva para nuevos retos y prioridades que se utilizará, entre otras cosas, para garantizar una respuesta adecuada de la Unión en caso de circunstancias imprevistas; hacer frente a nuevas necesidades o nuevos retos, como los surgidos en las fronteras de la UE o de sus vecinos, relacionados con situaciones de crisis y postcrisis o la presión migratoria; y promover nuevas prioridades o iniciativas internacionales o lideradas por la UE. Está reserva es uno de los elementos importados desde el FEO, y que ha permitido en el pasado hacer frente a circunstancias extraordinarias en un contexto altamente volátil como es el de las relaciones internacionales. Asimismo, el instrumento también prevé que los créditos de compromiso y de pago no utilizados sean automáticamente prorrogados al ejercicio siguiente, y que en caso de no utilización vuelvan a la línea presupuestaria de origen, a diferencia de las reglas estándar del presupuesto de la UE. La reserva está dotada con 9.534 millones de euros, un 1 1% del total.


Aplicación práctica  
Europa Global se limita a establecer, como sus predecesores, un cuadro normativo general y no una definición de cómo se debe gestionar concretamente la cooperación de la UE en cada una de las regiones, países socios o áreas de los programas temáticos cubiertos por dicho instrumento. Esta concreción, como era el caso de los anteriores instrumentos de cooperación, se lleva a cabo en dos etapas sucesivas:
• En primer lugar, la programac1on plurianual para establecer las prioridades y los montantes indicativos por región, país o tema a medio y largo plazo.
• En segundo lugar, la ejecución financiera de dichos programas plurianuales, que se materializa a través de decisiones de financiación que toman la forma de planes de acción y medidas individuales o especiales.

Sobre la primera etapa, la UE define sus ámbitos de actuación y objetivos específicos con cada país y región socios a través del proceso de programación. Los planes indicativos plurianuales (PlP) esbozan la base general de la cooperación de la UE en los países y regiones, y establecen la aplicación de la programación así como las prioridades, objetivos y resultados esperados, y los montantes indicativos, para la cooperación a medio y largo plazo con las regiones o países socios (programas geográficos) o en un área temática (programas temáticos). Estos documentos se redactan tras completar un proceso de consulta y dialogo con los países socios y otros actores relevantes, de acuerdo con los principios de la AOD de apropiación y de predictibilidad. Por regla general estos PlP tienen una duración de siete años, ajustándose al periodo presupuestario actual de la UE, pero se prevé una revisión de los mismos a medio plazo o la posibilidad de que su duración se alinee con los ciclos y planes de desarrollo de dichos países.

Como en el caso de los instrumentos de financiación exterior precedentes, la adopción de cada uno de los PIP bajo Europa Global por parte de la Comisión Europea requiere previamente laopinión favorable de los EEMM reunidos en un comité creado por el nuevo instrumento, deacuerdo con el proceso de "comitología". A través de estos comités, los EEMM ejercen el controlde la competencia de la Comisión de gestionar los programas comunitarios y el presupuesto de la Unión. Igualmente, a través de este mismo proceso de comitología, el Parlamento Europeo y el Consejo ejercen un control sobre la Comisión limitado a la verificación de que esta última no viola sus competencias.

Sobre la segunda etapa, en las decisiones financieras se establece el tipo de operación o de intervención de la Unión en la región, país o tema. En dichas decisiones se indican, por ejemplo, las modalidades (subvenciones, asistencia técnica, apoyo presupuestario, garantías presupuestarias), el calendario previsto y los montantes de dichas operaciones. Sobre base de estas decisiones se lanzan los procedimientos de licitación para suministrar asistencia técnica a los países beneficiarios; las convocatorias de subvenciones para financiar acciones de la sociedad civil; se negocian los convenios de financiación con los países socios para el apoyo presupuestario; y/o se delegan tareas de ejecμción presupuestaria a las agencias de desarrollo de los EEMM o a organizaciones internacionales como las Naciones Unidas o el Banco Mundial, con las cuales la UE colabora y trabaja conjuntamente.

Por último; la UE quiere trabajar junto con otros donantes, y en particular con sus EEMM y otros organismos europeos (en especial BEi y EBRO), en lo que se llama "programación conjunta" o "Team Europe approach". El propósito de este nuevo enfoque es lograr un mayor impacto y eficacia de la ayuda europea a terceros países. Los recursos y modalidades de implementación de estos actores se pondrían al servicio de iniciativas comúnmente acordadas de alto valor estratégico y con gran capacidad de transformación. El origen de este nuevo enfoque del Team Europe para el periodo 2021-2027 se encuentra en la exitosa respuesta conjunta ante el desafío de la pandemia del Covid- I 9, en la que, bajo la iniciativa de la Comisión Europea, la UE, los EEMM y sus instituciones financieras lograron movilizar en pocas semanas enormes paquetes financieros.


Instrumentos financieros

Así, casi el 72% de la acción exterior de la Unión está ahora armonizada dentro del nuevo instrumento (más del 85% si se incluyen el Instrumento de Ayuda de Preadhesión, la Decisión de Asociación de Ultramar y el Instrumento de Seguridad Nuclear). Dos son los instrumentos financieros concretos del nuevo sistema:
• El Fondo Europeo de Desarrollo Sostenible Plus (FEOS+). Los objetivos de gasto de este instrumento son:
o al menos el 93% del total destinarse a la AOD,
o debe contribuir a que la UE alcance el objetivo de entre el O, 1 5% '/ 0,20% de la RNB de la UE destinada a los PMA en forma de AOD;
o el 20% de la AOD debe financiar medidas de inclusión social y Desarrollo humano,
o el 10% del total debe destinarse a financiar la gestión y gobemanza de la migración y desplazamientos forzados;
o al menos el 85% de las nuevas acciones que se apliquen bajo el Instrumento Europa Global debe tener como objetivo principal o significativo la igualdad de género, y al menos el 5% de esas acciones debe tener la igualdad de género y los derechos de la mujer y su empoderamiento como objetivo principal, y . o el 30% del total de GE debe dedicarse a objetivos climáticos.
• La Garantía de Acción Exterior, con la capacidad de cubrir riesgos a terceros por un valor máximo de 53.449 millones de euros, con el objetivo de promocionar las inversiones en terceros países a través de dichos instrumentos financieros novedosos. De este modo, con la creación de incentivos para movilizar recursos privados, se logra un efecto multiplicador de los fondos europeos (estimados en la posibilidad de movilizar hasta 1 1 euros por cada euro de la UE a través de las garantías) y se ataja la llamada brecha financiera de los Objetivos de Desarrollo Sostenible, pasando de billones a trillones ("from millions to trillions"). 

Ciertos instrumentos financieros quedan fuera de Europa Global: el instrumento asistencia humanitaria, el instrumento de Preadhesión Ill (IPA III), la Facilidad Europea de paz, la Facilidad Europea para refugiados en Turquía (FriT) y los tres Trust-Funds de la UE creados en el periodo2014-2020 y extendidos hasta finales 2021 para responder a las crisis migratorias y de refugiados.


\begin{specialblock}{Formas adicionales de financiación}
    Blending es un mecanismo o plataforma de financiación combinada definido como "un marco de cooperación establecido entre la Comisión y las instituciones de desarrollo u otras instituciones financieras públicas, con el fin de combinar ayudas no reembolsables y/o instrumentos financieros y/o garantías presupuestarias del presupuesto y formas reembolsables de apoyo procedentes de instituciones de desarrollo u otras instituciones financieras públicas, así como de instituciones financieras y de inversores del sector privado". Las operaciones de financiación combinada son una de las diez fonnas de financiación de la UE previstas por el Reglamento Financiero que pueden utilizarse en el marco del Instrumento Europa Global. Las operaciones combinadas son formas de financiación que combinan subvenciones no reembolsables y préstamos reembolsables. Con garantías y blending, se utilizan fondos públicos para poner en marcha o cubrir parte de los costes de un proyecto, y pennite canalizar inversiones privadas en acciones de desarrollo.\footnote{%
        Las plataformas o instrumentos de financiamiento mixto o combinado (\textit{blending}) de la UE son mecanismos que impulsan el financiamiento de proyectos de inversión en los países beneficiarios de la cooperación externa de la UE.
    } 
    El Plan de Inversiones Exteriores (PE!) es una iniciativa emblemática de la UE lanzada en 2017, que utiliza tanto la financiación mixta corno las garantías. Su objetivo es atraer más inversiones, en particular de empresas e inversores privados, a los países vecinos de la UE (vecinos de la UE) y a África. Se espera que movilice 47.000 millones de euros de inversión a través de una aportación de la UE de 4.600 millones de euros.
    
    2. Los Trust-Funds_ son instrumentos extrapresupuestarios con su propio sistema de financiación que combina fondos europeos y contribuciones de EEMM y otros donantes. Se vienen creando para reunir financiación de distintas fuentes de manera temporal. La UE puede crear fondos fiduciarios de varios donantes para acciones de emergencia, post-emergencia o temáticas en el ámbito de la acción exterior. Desde 2014, se han creado cuatro fondos fiduciarios de la UE: el Fondo Fiduciario de la UE para Colombia que no se va a renovar a partir de 2022, el Fondo Fiduciario de Emergencia de la UE para la estabilidad y para abordar las causas profundas de la migración irregular y los desplazados en África (Fondo Fiduciario de la UE para África), el Fondo Fiduciario de la UE en respuesta a la crisis siria (Fondo Fiduciario "Madad" con 2 200 millones de euros) y el Fondo Fiduciario de la UE para la República Centroafricana (Fondo Fiduciario "Bekou" (1 08 millones de euros).
\end{specialblock}

\paragraph{Instrumentos de financiación (\textit{ex} art. 209.3): El papel del Banco Europeo de Inversiones}
Entre los instrumentos de la cooperación hay que mencionar la actividad del Banco Europeo de Inversiones (BEI) que, alentado por la flexibilidad de sus Estatutos en este campo (art. 209 TFUE): [El BEI] \textit{contribuirá, en las condiciones previstas en sus Estatutos, a la ejecución de las acciones contempladas en el apartado 1} para el logro de los objetivos de la política de cooperación al desarrollo.

Dentro del entramado de instituciones que llevan a cabo la implementación de la política de financiación al desarrollo de la Unión Europea hay que mencionar al Banco Europeo de Inversiones (BEi). El BEi es una institución financiera multilateral, cuyos socios son los 27 países de la UE. Como en otras instituciones multilaterales de desarrollo, el capital social del BEi proviene de las aportaciones de sus socios y además el BEi se endeuda en los mercados financieros en los que goza de un rating AAA, para obtener recursos para prestarlos, a tipos y plazos competitivos para proyectos de inversión, tanto del sector público como del sector privado 

Aunque la mayor parte de la financiación del BEi se concede dentro de la UE (en 2022 se firmaron inversiones dentro de la UE por valor de unos 57.000M€ frente a unos 8.800M€ fuera de la UE, es decir la actividad fuera de la UE supone solo alrededor del 13\% de la actividad del BEI), el gran tamaño del banco hace que su actividad fuera de la UE sea comparable con otros bancos de desarrollo (a modo de ejemplo el BERD aprobó en 2021 préstamos por un valor de alrededor de I O.OOOM€, pero una parte es dentro de la UE o el Banco Interamericano de Desarrollo aprueba anualmente proyectos por valor de 14.000M$).

El Consejo Europeo de junio de 2021 , en sus Conclusiones sobre el refuerzo de la arquitectura financiera europea para el desarrollo invitó al BEJ a que "presente medidas de mejora que refuercen la repercusión en materia de desarrollo de sus operaciones en los países socios". El 1 de enero de 2022, como respuesta a esta petición, se creó "EIB Global", una sucursal del BEi con su propia marca e identidad que gestionará los proyectos fuera de la UE.

Cuando el BEi actúa fuera de la UE, tradicionalmente lo hace con apoyo de la UE en forma de garantías o subvenciones a los tipos de interés, ya que su prudente política de riesgos, si no, le limitaría a la hora de invertir en los países más frágiles. Así el BEi será un agente muy importante en la ejecución del nuevo Fondo Europeo de Desarrollo Sostenible Plus (FEDS+). El BEi y la Comisión firmaron un acuerdo de garantía que permitirá a la UE cubrir el riesgo de los préstamos del BEi fuera de la UE y de esa forma permitir al Banco financiar proyectos en apoyo de la política de la UE en entornos de mayor riesgo. Sin embargo, en los últimos años y dada su buena posición de capital, ha comenzado a asumir él mismo el riesgo de algunos de en sus préstamos fuera de laUE (aunque siguen siendo una minoría).

La mayoría de los proyectos que apoya se cofinancian con otras instituciones de financiación del desarrollo, como el Banco Mundial, el Banco Europeo de Reconstrucción y Desarrollo (BERD), el banco de desarrollo alemán KtW, la Agence Française de Développement, el Banco Asiático de Desarrollo, el Banco Africano de Desarrollo y el Banco Interamericano de Desarrollo, entre otros. Estos proyectos pretenden desarrollar la industria, la innovación y las infraestructuras, con el objetivo último de corregir los desequilibrios y promover una economía estable. Los proyectos del BEi crean empleo, generan energía más barata y limpia y propician una mayor igualdad de género e integración regional.


\subsubsection{Naturaleza convencional (\textit{ex} art. $209.2$): Acuerdos de Asociación Económica (EPA, ACP)}
Los acuerdos internacionales en materia de cooperación al desarrollo (art. $209.2$) se celebrarán de conformidad con lo dispuesto en el art. $218$ TFUE, donde se recoge el procedimiento general de celebración de acuerdos, que supondrá la participación del Parlamento Europeo mediante la aprobación o la consulta. Hoy en día adoptan la forma de Acuerdos de Asociación Económica (EPA, por sus siglas en inglés).\footnote{%
    Estos constituyen la respuesta de la Unión a un problema jurídico concreto planteado en el seno de la OMC. Las preferencias comerciales no recíprocas que la Unión concedía a los países de África, el Caribe y el Pacífico bajo los sucesivos Convenios de Yaundé (1963-1975) y Lomé (1975-2000) eran, en principio, incompatibles con la cláusula de nación más favorecida del artículo I del GATT, al beneficiar selectivamente a un grupo de países en desarrollo y no a la totalidad de estos, lo que impedía ampararlas en la cláusula de habilitación (\textit{enabling clause}), aplicable únicamente a esquemas de preferencias universales como el SPG. Su mantenimiento solo fue posible gracias a sucesivas exenciones (\textit{waivers}) del art. 1 del GATT concedidas por la OMC, la última de ellas obtenida en la Conferencia Ministerial de Doha de 2001, con vencimiento el 31 de diciembre de 2007. Sobre la incompatibilidad de las preferencias de Lomé con el GATT y el mecanismo del \textit{waiver}, véase FAO, "La situación de las preferencias comerciales en la OMC"; y el litigio \textit{Comunidades Europeas -- Régimen para la importación, venta y distribución de bananos} (DS27), donde el Órgano de Apelación de la OMC constató la pérdida de cobertura de la exención de Doha aplicada a los bananos ACP a partir del 1 de enero de 2006.
} El Acuerdo de Cotonú (2000), que sucedió a Lomé, prorrogó las preferencias no recíprocas únicamente durante un período preparatorio (2000-2007), con el mandato expreso de sustituirlas, antes de esa fecha, por acuerdos comerciales recíprocos y compatibles con las normas de la OMC: los EPA.

Aunque no constituye una agrupación política, ni un bloque comercial, el grupo ACP —África, Caribe y Pacífico— ha sido el principal beneficiario de la cooperación al desarrollo de la Unión Europea. La condición de antiguas colonias de sus países lo situó en un lugar privilegiado en las relaciones exteriores de la Comunidad. De hecho, las convenciones firmadas entre ambas partes conformaron un sistema, conocido habitualmente como \textbf{Sistema Lomé}, que ha sido considerado un modelo único en las relaciones Norte-Sur, ya que es su mayor marco político, comercial y financiero. Aunque, de los dos pilares en los que se ha sustentado la amplia cooperación UE-ACP, la ayuda —financiera o técnica— y el comercio, ha sido este último el que haalcanzado un mayor protagonismo. 

El marco político general que envuelve actualmente los EPA es el Acuerdo de Samoa, firmado el 15 de noviembre de 2023 entre la Unión y los $79$ miembros de la Organización de Estados de África, el Caribe y el Pacífico (OEACP). Constituye un marco común aplicable a los $79$ Estados, complementado por tres protocolos regionales específicos para África, el Caribe y el Pacífico, y regirá la relación política, de cooperación y de derechos humanos durante los próximos veinte años; no sustituye, sin embargo, a los EPA regionales, que permanecen en vigor como pilar estrictamente comercial de la relación, ahora enmarcado bajo el paraguas político del nuevo Acuerdo.

\paragraph{Protocolo Caribe: AAE CARIFORUM-UE}
El Acuerdo de Asociación Económica (2008) con los catorce Estados del CARIFORUM (los miembros de la Comunidad del Caribe más República Dominicana), fue el primer EPA celebrado por la Unión, y se ha convertido en el modelo de referencia para las negociaciones posteriores con las demás regiones ACP.

A diferencia del resto de EPA, que en su mayoría se han limitado durante años al comercio de mercancías, este incorpora desde su origen compromisos sustanciales en materia de comercio de servicios (con los Estados más desarrollados del CARIFORUM comprometiéndose a liberalizar el $75\%$ de su comercio de servicios con la Unión, y los menos desarrollados el $65\%$, inversiones, política de competencia, contratación pública y propiedad intelectual, además de un Protocolo específico de cooperación cultural. Pese a esta profundidad de contenido, su fundamento último sigue siendo asimétrico: el acceso de la Unión es inmediato, mientras que los Estados del CARIFORUM disponen de hasta veinticinco años para liberalizar sus importaciones procedentes de la Unión.

\paragraph{Protocolo África: Comunidad para el Desarrollo del África Meridional (SADC, 2016)}
El AAE UE-SADC, firmado el 10 de junio de 2016 por Botsuana, Esuatini, Lesoto, Mozambique, Namibia y Sudáfrica, se aplica provisionalmente desde el 10 de octubre de 2016. Este acuerdo reviste una importancia particular por haber sustituido, en su vertiente estrictamente comercial, al Acuerdo de Comercio, Desarrollo y Cooperación (TDCA) que la Unión mantenía bilateralmente con Sudáfrica desde 1999.\footnote{%
    Sudáfrica ocupa dentro del EPA de la SADC una posición singular y menos favorable que el resto de socios. Por un lado, mientras que Botsuana, Esuatini, Lesoto, Mozambique y Namibia (de la Unión Aduanera del África Austral, SACU) reciben de la Unión un acceso libre de aranceles y contingentes para prácticamente la totalidad de sus exportaciones (salvo armas y municiones), Sudáfrica, por su mayor nivel de desarrollo relativo, solo obtiene ese trato para el $96\%$ de sus exportaciones, beneficiándose del restante únicamente de una reducción arancelaria o de contingentes preferentes. En sentido inverso, el conjunto SACU abre su mercado a la Unión de forma libre de aranceles y contingentes solo para el $84,9\%$ de las importaciones europeas.
}

No obstante, frente al caso relativamente ordenado del SADC, el resto de configuraciones regionales africanas del EPA ilustran, más que un régimen consolidado, un proceso de implementación fragmentado y, en buena parte, todavía inconcluso. Por un lado, el AAE con los dieciséis Estados de África Occidental (CEDEAO) negociado como bloque regional, no se aplica en la actualidad en ninguno de los países de la región. La razón es que su entrada en aplicación provisional exige que la totalidad de los Estados lo hayan firmado y que al menos dos tercios lo hayan ratificado, condición que continúa sin cumplirse porque Nigeria no lo ha firmado, alegando el riesgo de desindustrialización que supondría para su sector manufacturero la apertura recíproca frente a la Unión.

En África Central únicamente Camerún aplica, desde 2014, un EPA interino y limitado al comercio de mercancías, mientras que el resto de Estados de la región permanecen fuera del acuerdo y se benefician, en su lugar, del régimen unilateral del SPG o, en el caso de los países menos desarrollados, del régimen EBA.

Las negociaciones del AAE con la Comunidad del África Oriental (EAC) concluyeron en octubre de 2014, y el acuerdo ha sido firmado por Kenia, por Ruanda y por todos los Estados miembros de la Unión, pero continúa sin entrar en vigor a la espera de la firma del resto de socios de la EAC. Este bloqueo resulta especialmente ilustrativo porque, al tratarse de un acuerdo con una unión aduanera regional ya constituida (el Arancel Exterior Común de la EAC), su falta de entrada en vigor genera además fricciones internas dentro del propio bloque regional entre los Estados que sí querrían aplicarlo --Kenia, principal exportador de la región a la Unión y único Estado de la EAC no clasificado como país menos desarrollado, por lo que sin el EPA perdería su acceso preferente al mercado europeo-- y los que se resisten a hacerlo.

Esta implementación desigual es señalada de forma recurrente por la doctrina como evidencia de que el formato de negociación por bloques regionales, pensado para reforzar la integración económica sur-sur, ha topado en la práctica con las asimetrías de intereses económicos entre los propios Estados de cada bloque, retrasando los beneficios que el propio EPA pretendía generar para el conjunto de la región.

\paragraph{Protocolo Pacífico: Estados del Pacífico}
El AAE interino con los Estados ACP del Pacífico fue firmado inicialmente por Papúa Nueva Guinea (julio de 2009) y por Fiyi (diciembre de 2009), aplicándose provisionalmente Papúa Nueva Guinea tras su ratificación en 2011 y Fiyi desde 2014; posteriormente se han adherido Samoa (diciembre de 2018) y las Islas Salomón (mayo de 2020). De los catorce Estados ACP de la región, Papúa Nueva Guinea y Fiyi concentran, con diferencia, el grueso del comercio bilateral con la Unión, lo que explica que el resto de Estados insulares del Pacífico --de reducidísimo peso comercial-- no hayan mostrado el mismo interés en completar su adhesión al Acuerdo.






\subsubsection{La financiación}
El Capítulo 1 del Título III del TFUE no contiene una disposición general sobre la financiación de la política de cooperación al desarrollo. Sin embargo, el contenido del artículo permite extraer algunas conclusiones en materia de financiación:
\begin{la}
    \item \textbf{Sistema doble de financiación}. En primer lugar, el mantenimiento de un sistema de doble financiación, el de la Unión y el nacional, de las actividades de cooperación al desarrollo. Esto se explica por varias razones: la articulación jurídica misma de esta política que tiene un carácter complementario respecto de las políticas nacionales, la previsión de una coordinación y de una concertación entre los Estados miembros y la UE en sus políticas y programas de ayuda y la posibilidad de realizar acciones conjuntas y de que los Estados contribuyan, \textit{si fuere necesario, a la ejecución de los programas de ayuda de la Unión}; y,
    \item \textbf{Participación del BEI}. En segundo lugar, la previsión contenida relativa al Banco Europeo de Inversiones, que legitima su actuación, con carácter general, en la ejecución de las acciones de cooperación al desarrollo mediante la financiación de los correspondientes proyectos de cooperación, como ya se venía haciendo, con anterioridad, en relación con los países ACP y los mediterráneos.
\end{la}

En términos generales, se confirma la práctica anterior en materia de financiación, basada en la utilización privilegiada del Fondo Europeo de Desarrollo.\footnote{%
    La prueba más evidente de esa continuidad se encuentra en el hecho de que en una Declaración anexa al Tratado de Maastricht relativa al FED se «conviene en que el Fondo Europeo de Desarrollo seguirá financiándose mediante contribuciones nacionales, de conformidad con las disposiciones vigentes». Con ello se hizo caso omiso a las pretensiones del Parlamento Europeo y de la Comisión que propugnan la inclusión del FED en el presupuesto de la Unión. Con esa inclusión se conseguiría asociar las Instituciones de la Unión al control de este instrumento financiero principal de la cooperación al desarrollo, que de este modo sigue estando sujeto a las discusiones entre los Estados miembros motivadas por la fórmula del reparto de los gastos.
}

\subsection{Ayuda Oficial al Desarrollo (AOD)}
Siguiendo el Comité de Ayuda al Desarrollo (CAD) de la OCDE, la AOD se define corno el conjunto de recursos dirigidos a los PED con el objetivo de fomentar su progreso socioeconórnico, y que cumplen unos requisitos de concesionalidad establecidos por el propio CAD (esencialmente, contar con un componente de donación de, al menos, un 25\% del total de los recursos transferidos).\footnote{%
    El CAD está compuesto actualmente por 30 miembros: Australia, Canadá, Islandia, Japón, Corea, Nueva Zelanda, Noruega, Suiza, Reino Unido y Estados Unidos, $19$ Estados miembros de la UE y la Unión Europea que actúa como miembro de pleno derecho del comité, 3 Estados miembros de la UE (Estonia. Letonia y Lituania) no son miembros del CAD de la OCDE. mientras que otros 5 (Bulgaria Croacia. Chipre. Malta y Rumanía) no son miembros ni de la OCDE ni del CAD. pero Bulgaria y Rumanía sí son  participantes en el CAD. El CAD se dedica a coordinar a los donantes. y establece una reglas mínimas para garantizar la eficiencia de la AOD.
}




El CAD revisa la lista de países elegibles para la AOD cada tres años. La lista del CAD de receptores de AOD muestra todos los países y territorios que pueden recibir AOD. Se trata de todos los países de renta baja y media según la renta nacional bruta (RNB) per cápita publicada por el Banco Mundial, a excepción de los miembros del G8, los miembros de la UE y los países con fecha firme de ingreso en la UE. La lista también incluye a todos los PMA según la definición de las Naciones Unidas.


1. Cooperación financiera
Supone la asignación de recursos financieros a los PEO, de manera bilateral o multilateral (a través de instituciones financieras multilaterales, como el BM). A su vez, la cooperación financiera puede ser no reembolsable, es decir, sin reintegro de los fondos otorgados; o reembolsable, que asume la forma de créditos blandos, en condiciones de tipos de interés y plazos más favorables que los del mercado.

2. La asistencia técnica 
Implica el apoyo a los PEO mediante la transferencia de conocimientos, tecnologías, técnicas productivas, o habilidades y experiencias. La cooperación financiera y la asistencia técnica es la forma en que la UE (y el resto de los países desarrollados) ha llevado a cabo los proyectos tradicionales de cooperación al desarrollo, como puede ser, por ejemplo, la construcción de un hospital. Otro componente clásico de estas dos modalidades de AOO es la ayuda alimentaria normal (no de urgencia), que, históricamente,consistía en el envío de alimentos, y que, hoy, fomenta una mayor seguridad del suministro alimentario en de los países beneficiarios. Otros componentes más modernos de estas formas de AOO pueden ser la financiación de proyectos realizados por las ONG.


3. La cooperación económica

Está destinada a promover proyectos y actividades de interés común para los agentes económicos y sociales de los países donantes (países en desarrollo) y los países receptores (PEO). Se incluyen aquí, por ejemplo, las acciones que tienden a facilitar el comercio (Aid for Trade).\footnote{%
    La estrategia de_Ayuda al Comercio de la UE se adoptó en 2007 en respuesta a la iniciativa de Ayuda al Comercio que la Organización Mundial del Comercio (OMC) lanzó en 2005.
} o a fomentar la inversión extranjera en el país receptor.

En relación a las ventajas comerciales con ciertos socios (este tema es abordado en profundidad en el tema 3A4 I ), el actual Sistema de Preferencias Generalizadas (SPG) de la UE viene regulado en el Reglamento (UE) 978/201 y se aplica desde el I de enero de 2014. Los tres regímenes del sistema, el SPG general, el SPG+ de incentivos y el sistema "Todo menos armas" (TMA) se refuerzan ajustando las preferencias y garantizando que tengan un mayor impacto. Los países que pueden acogerse al SPG figuran en el anexo I del Reglamento SPG. Los países beneficiarios de las nuevas preferencias SPG figuran en el Anexo 11. Los beneficiarios de "Todo menos armas" figuran en el anexo IV. Para beneficiarse de este tratamiento preferencial, los  productos deben cumplir con las reglas de origen, esto es, proceder de los países incluidos en los anexos anteriormente mencionados.

En 2020, coincidiendo con el inicio de la pandemia, que tuvo un impacto dramático en el comercio, la UE y sus Estados miembros invirtieron colectivamente 23.000 millones de euros en facilitar el comercio (Aid for Trade). La mayor parte se destinó a África, seguida de Asia. En junio de 2022, 123 países y territorios socios podían optar a la AOD de la UE y tenían acceso preferencial al mercado de la UE: 62 de ellos a través de un acuerdo comercial preferencial en vigor, 61 a través de uno de los tres regímenes comerciales preferenciales unilaterales de la UE en el marco del SPG. La estrategia de la UE en materia de ayuda al comercio está concebida para prestar una ayuda más pertinente que haga uso de todos los instrumentos de la UE, combinando la AOD con incentivos a la inversión y el acceso a los mercados libre de derechos ofrecidos a través de los acuerdos de libre comercio y los regímenes preferenciales. La UE se propone seguir aumentando la parte de la ayuda al comercio asignada a sus países socios para ayudarles a duplicar su cuota de exportaciones mundiales de aquí a 2020- 2025.

4. La ayuda humanitaria y de emergencia

Tiene como objeto prestar asistencia y socorro a PMA y poblaciones con grandes necesidades derivadas de catástrofes naturales o de origen humano (como los conflictos armados). Dicha asistencia incluye la ayuda alimentaria de urgencia, la asistencia a los refugiados y repatriados, el envío de bienes y equipos de primera necesidad, y la atención de otras necesidades humanitarias derivadas de tales catástrofes.

Situación actual e11 la VE 
El informe anual del Consejo sobre los objetivos de ayuda al desarrollo señala que en 202 1 la UE mantuvo su posición como mayor proveedor mundial de AOD habiendo destinado 70.200 millones de euros, un $4,3$\% más que en 2020.\footnote{%
    De esos 70.200 millones de dólares. aproximadamente. unos 855 millones provienen de la U E como institución ( 1 .3\%). y los restantes 69.345 millones (98.7\%) son ayuda bilateral de los Estados miembros. Estos datos reflejan la estructura previa a Europa Global.
} La UE y sus Estados miembros, trabajando juntos desde el Team Europe approach, confirman así su posición como primer donante mundial, aportando el 43\% de la AOD mundial. El compromiso que la UE adquirió en 2015 en la Agenda de Acción Adís Abeba sobre la Financiación para el Desarrollo (que es parte de la Agenda 2030) es

• Aportar al menos el 0,7\% de la RNB colectiva como AOD para 2030 ... Actualmente la UE destina el 0,49\% de la RNB de la UE frente al 0,26\% del conjunto de donantes del CAD no pertenecientes a la UE y solo cuatro EEMM ya han alcanzado o superado ese objetivo en 2021· : Luxemburgo (0,99\%), Suecia (0,92\%), Alemania (0,74\%) y Dinamarca (0,70\%) (España el 0,25\%).\footnote{%
    Tras la caída de la AOD como consecuencia de las restricciones presupuestarias nacionales y comunitarias impuestas a raíz de la crisis financiera económica internacional, estos flujos se fueron recup��rando hasta alcanzar un máximo en 2016 (0.52% de la RNB mientras que EE.UU. destinó el 0, 1 9%, y Japón el 0.20%. Posteriormente. ese porcentaje se redujo gradualmente hasta 20 19. En 2020 y 2021 la tendencia es creciente de nuevo. (UE: 0.49%. EE.UU.: 0. 1 8%. Japón: 0,34%).
}


• ... con un componente para los países menos desarrollados del O, 1 5\% al 0,20\% de la RNB a corto plazo y el 0,20% para 2030 . En 2020, la AOD colectiva de la UE y sus 27 EMM a los PMA alcanzó los 1 6.300 millones de euros, es decir, el 0, 12% de la RNB. Sólo cinco EEMM alcanzaron o superaron este umbral del O, 1 5% de AOD a los PMA/INB en 2020 : Luxemburgo (0 ,48%), Suecia (0,35%), Dinamarca (0 ,21 %), Bélgica (O,1 6%) y Alemania (O, 15%).

A nivel geográfico, la AOD de la UE financia programas de desarrollo en unos 140 países en todo el mundo, además de los programas de carácter regional o global. Casi el 50% de la AOD de la UE se dirige al África Subsahariaña, lo que refleja la especial preocupación que la UE sigue mostrando por la región, que constituye, en su conjunto, la zona más pobre de la Tierra, e, históricamente, ha sido una zona de presencia europea. Más en detalle, se puede observar que las principales regiones beneficiarias de la AOD europea en 2020 eran África Subsahariana (23%), Asia central y sureste asiático (9%), Europa (8,2%), Oriente Medio (6,5%):\footnote{
A pesar de que la región de África Subsahariana es la principal destinataria de la AOD. la ayuda de la UE se ha caracterizado por una creciente diversificación geográfica en el tiempo, que contrasta con el caso de EEUU o Japón, cuyo apoyo al desarrollo está más concentrado en determinados países. Hoy en día, la distribución de la AOD de la UE es más uniforme de lo que era hace algunas décadas, reflejando la conciencia de actor global que mantiene la Comisión Europea.
}

En relación a la distribución sectorial de la AOD de la UE, ·en este caso, únicamente nos centraremos en la procedente de las instituciones de la UE (obviando la de los EEMM). Para este análisis, nos basamos en la clasificación sectorial que realiza el CAD de la OCDE. Los principales sectores destinatarios de la AOD bilateral comprometida de las instituciones de la UE fueron en 2020 Gobierno y sociedad civil (1 9,6%), educación ( 12,5%), respuesta de emergencia ( 1 2,2%), ayuda a refugiados que creció a raíz de la crisis en Siria (9,5%), y salud (9, 1 %):


Como se puede observar, l a U E apuesta por estos sectores esenciales: infraestructura social, ayuda humanitaria, multisectorial y sectores de producción. Como en otras materias, la Comisión gestiona gran parte de los fondos destinados a la cooperación al desarrollo.

En línea con el enfoque actual sobre el desarrollo, la UE busca contribuir a crear el entorno adecuado, con estabilidad política y social, instituciones sólidas, regulación favorable al mercado, infraestructuras para propiciar el despegue de los países beneficiarios y que la AOD sea lo más efectiva posible. Todo lo anterior, siguiendo las directrices internacionales, y coordinando a los Estados miembros para asegurar la eficiencia de la cooperación al desarrollo.


\begin{specialblock}{TEI y Global Gateway}
    La AOD gestionada por la Comisión Europea ascendió a 1 5 . 1 00 millones de euros, incluidos 1 .000 millones de pré_stamos especiales de ayuda macrofinanciera en régimen equivalente de subvención, algo menos que en 2020. La AOD gestionada por el BEi ascendió a 1 .000 millones de euros.
    
    Como se mencionó en el anterior apartado, Team Europe approach, consistente en la colaboración entre la UE, los EEMM y el BEi y EBRO tiene aunar por fin recursos y mejorar la coordinación y visibi lidad para conseguir el mayor impacto sostenible aprovechando la riqueza de la experiencia colectiva, los conocimientos especializados y los recursos. Las Iniciativas Team Europe (TE!) son el buque insignia del enfoque Team Europe. El análisis conjunto que conduce a la identificación de una TEI a nivel nacional debe identificar una prioridad crítica que esté limitando el desarrollo y en la que un esfuerzo coordinado y coherente por parte de Team Europe garantizaría un impacto transformador. Las TEI también deben reunir la mejor combinación posible de modalidades, herramientas y socios para lograr el impacto previsto. Las TEI se basan en las prioridades políticas y normativas de la UE y se financian tanto con cargo al presupuesto de la U E -y, por lo tanto, se rigen por las directrices de programación de la NDICI-, como con cargo a los recursos financieros pertinentes de los Estados miembros de la UE participantes. Los procesos de programación conjunta a nivel nacional deben reflejar e incorporar las TE! y, a su     vez, las TE I que se identifiquen deben alimentar los procesos de programación conjunta en curso y futuros. También puede haber TEI plurinacionales o regionales, cuando puedan garantizarse las sinergias y la eficacia a este nivel . Capacity4dev es la plataforma de la Comisión dedicada a la difusión del conocimiento sobre Cooperación internacional y Desarrollo e infonna de las múltiples TEI que se están desarrollando (en total son 169, de las cuales 133 están específicamente asignadas a países, siendo 32 regionales y 4 globales). 

    Asimismo, la web de Team Europe Initiative and Joint Programming tracker presenta información detallada de las TEI y los procesos de programación conjunta que pueden servir al opositor para tener una idea más precisa de los distintos tipos de proyectos por regiones, países de origen y destino, sectores. Por ejemplo, España participa en 56 TEI (42,1 % del total) y en 1 9 TEI regionales (59,4% del total).

    En 2021, en el marco del Team Europe approach, se puso en marcha la estrategia Global Gateway para movilizar inversiones por un total de hasta 300.000 millones de euros durante el período 20 21 -2027. En consonancia con las ambiciones geopolíticas de la UE y el compromiso con la Agenda 20 30, la estrategia Global Gateway tiene por objeto impulsar vínculos inteligentes,limpios y seguros en los ámbitos digital, energético y del transporte, y reforzar los sistemas de salud, educación e investigación en todo el mundo para sustentar una recuperación mundial duradera, promoviendo al mismo tiempo valores universales y normas estrictas, buena gobemanza y transparencia
\end{specialblock}


Valoración
A pesar de los importantes avances expuestos más arriba, siguen persistiendo problemas de coordinación de las políticas de cooperación nacionales con la europea. Además, los objetivos puros de desarrollo pueden acabar chocando con otras medidas en los ámbitos comercial, migratorio, o de seguridad. Todo ello puede acabar menoscabando la coherencia y eficacia de las acciones de desarrollo. Por ello son necesarios informes periódicos que examinen el impacto real de la política de cooperación al desarrollo. Por otro lado, algunos expertos consideran que el dialogo político debería ser más inclusivo permitiendo que más actores (parlamentos, sindicatos, y partidos políticos) tuvieran voz en estos debates. Según su visión, ello mejoraría el apoyo social y también presupuestario de las medidas aplicadas y por tanto su sostenibilidad. Otros consideran incluso necesario incluir salvaguardias que protejan los fondos europeos de acabar fortaleciendo regímenes poco democráticos. Asimismo, otro de los retos a los que se enfrentan las delegaciones de la UE en el mundo encargadas de aplicar la política de cooperación al desarrollo a nivel localson las restricciones organizativas y de personal.










































































\clearpage
\section{Conclusión} 


\subsection{Recapitulación}


\subsection{Extensiones y relación con otras partes del Temario}



\subsection{Opinión}

\subsubsection{Idea final (Salida o cierra)}

\clearpage
\pagenumbering{gobble}
\MyBibliography



\MyBibItem{DAWID y GATTI}{Chapter 2 - Agent-Based Macroeconomics}{2018}{https://doi.org/10.1016/bs.hescom.2018.02.006}


% ANEXOS
\clearpage
\anexo{\textbf{Preguntas Test}}
%\input{\TemaCode/questions.tex} 







\clearpage
\anexo{Políticas Europeas de Adhesión y Vecindad}\label{anx:politica_vecindad}

\href{https://www.europarl.europa.eu/factsheets/es/sheet/170/la-politica-europea-de-vecindad}{Política Europea de vecindad}












\end{document}